%% =====================================================================
%%  Capítulo 6 -- Bifurcaciones
%%  Formato nuevo (noviembre de 2026). Integra: Handout5.tex (versión anterior:
%%  estabilidad estructural, bifurcaciones 1-D y 2-D, Hopf, "blue sky" y mapas
%%  de Poincaré), las diapositivas Bifurcaciones2026 y sus complementos
%%  pedagógicos (Bifurcaciones2D_Complementos: hiperbolicidad y las dos rutas,
%%  receta de cinco pasos, el fantasma, el interruptor genético, genericidad,
%%  Hopf completo, histéresis de ciclos, preguntas discriminantes) y los
%%  scripts bif_*.m, BifEj*.m y Bifurcaciones80min_LiveScript.m, reescritos
%%  en codigo/cap06_bifurcaciones.
%%  Se agregaron: la condición de la función implícita, las condiciones de
%%  Sotomayor, la horquilla escondida (variedad central con parámetro), el
%%  coeficiente de Lyapunov, el retardo de la bifurcación con rampa lenta,
%%  la fórmula de la divergencia para el multiplicador y el mapa de Poincaré
%%  numérico de van der Pol. Correcciones: ver \S\ref{sec:bf-cierre}.
%%  Números verificados con herramientas/cap06_numeros.py (125/125) y
%%  codigo/cap06_bifurcaciones/verificarCap06.m (42/42 en Octave 8.4).
%% =====================================================================

\capitulo[cap:bifurcaciones]{Bifurcaciones}
  {Otra consecuencia es que se está impulsando que la educación sea vista como una capacitación que solo conduce al trabajo, y no como un medio que conduce al pensamiento crítico y a la transformación.}
  {Junot Díaz}

\begin{objetivos}[partes/06-bifurcaciones.html]
  \item Explicar qué es una bifurcación local como pérdida de hiperbolicidad, y reconocer en el plano traza--determinante las dos maneras en que puede ocurrir en $\R^2$.
  \item Analizar las bifurcaciones silla-nodo, transcrítica y de horquilla (supercrítica y subcrítica) con líneas de fase y diagramas de bifurcación, y decidir cuáles sobreviven a una perturbación del modelo.
  \item Seguir una receta de cinco pasos (equilibrios, jacobiano, cruce, transversalidad y no degeneración) para encontrar y clasificar bifurcaciones en sistemas planos, reduciendo a la variedad central cuando haga falta.
  \item Enunciar el teorema de Hopf con todas sus hipótesis, distinguir la bifurcación supercrítica de la subcrítica con el primer coeficiente de Lyapunov o con una simulación, y explicar la histéresis que produce una subcrítica.
  \item Construir el mapa de Poincaré de una órbita periódica, calcular su multiplicador (analíticamente o con la divergencia) y decidir con él la estabilidad del ciclo.
\end{objetivos}

%% ---------------------------------------------------------------------
\section{¿Qué cambia cuando cambia un parámetro?}
\label{sec:bf-porque}
%% ---------------------------------------------------------------------

Hasta aquí los parámetros de cada modelo estaban fijos, y la pregunta era qué hacen las trayectorias. Pero en ingeniería los parámetros se mueven: la temperatura cambia la característica de un resistor, una carga crece sobre una columna, un gen se degrada más rápido, la potencia que se bombea a un láser sube. Casi siempre un cambio pequeño de parámetro produce un cambio pequeño en el retrato de fase: los equilibrios se corren un poco y conservan su tipo. Este capítulo trata de las excepciones, los valores del parámetro donde un cambio arbitrariamente pequeño cambia la \emph{forma} del retrato: aparecen o desaparecen equilibrios, uno estable se vuelve inestable, nace una oscilación.

El ejemplo clásico es una columna vertical con un peso encima \citep[\S3.0]{Str15}. Con poco peso la columna está recta, y si se la empuja vuelve a quedar recta. Al aumentar el peso nada cambia\dots{} hasta una carga crítica, a partir de la cual la posición recta deja de ser estable y la columna se pandea hacia un lado o hacia el otro. Nada en el modelo cambió de manera abrupta; solo un parámetro cruzó un valor. A ese fenómeno lo llamaremos bifurcación (de horquilla, en este caso, porque el estado recto se divide en dos ramas simétricas).

Consideramos familias de sistemas
\begin{equation}
  \dot x=f(x,\mu),\qquad x\in\R^n,\quad \mu\in\R ,
  \label{eq:bf-familia}
\end{equation}
\emph{invariantes en el tiempo} y con $f$ suave en $(x,\mu)$; $\mu$ es el \clave{parámetro de bifurcación}. Para decir con precisión qué significa ``cambiar la forma del retrato'' necesitamos comparar retratos.

\begin{definicion}[label=def:bf-equivalencia]{Equivalencia topológica y estabilidad estructural}
Dos sistemas $\dot x=f(x)$ y $\dot x=g(x)$ son \clave{topológicamente equivalentes} en una región si hay un homeomorfismo (una deformación continua con inversa continua) que lleva las trayectorias de uno a las del otro conservando el sentido del tiempo. Un campo $f$ es \clave{estructuralmente estable} si todo campo $g$ suficientemente cercano a $f$ (cerca en $f$ y en sus primeras derivadas) es topológicamente equivalente a $f$.
\end{definicion}

\notamargen{Es la misma noción de ``igual forma'' del teorema de Hartman--Grobman (\cref{thm:bi-hg}), pero ahora global y aplicada a la familia de todos los campos cercanos, no a la linealización.}

El oscilador de Lotka--Volterra (\cref{sec:oc-lotka}) es el ejemplo típico de un campo estructuralmente \emph{inestable}: sus órbitas cerradas forman una familia continua alrededor de un centro, y cualquier perturbación que agregue un poco de fricción o de competencia intraespecífica las destruye. Andronov y Pontryagin caracterizaron en 1937 los campos estructuralmente estables en el disco, y Peixoto extendió el resultado a superficies compactas en 1962~\citep{AnPo37,Pei62}.

\begin{teorema}[label=thm:bf-peixoto]{Andronov--Pontryagin \citep[\S4.3]{Per01}}
Sea $f\in\Cont1$ en una región $D\subset\R^2$ cerrada y acotada, cuya frontera es una curva suave que el flujo cruza transversalmente (hacia adentro). Entonces $f$ es estructuralmente estable en $D$ si y solo si
\begin{enumerate}
  \item tiene un número finito de equilibrios y de órbitas periódicas, y todos son \alerta{hiperbólicos};
  \item no hay trayectorias que conecten dos sillas (ni una silla consigo misma).
\end{enumerate}
\end{teorema}

\notamargen{En el disco no hace falta una tercera condición que sí aparece en el teorema de Peixoto para superficies (que los únicos conjuntos límite sean equilibrios y órbitas periódicas): en el plano la garantiza el teorema de Poincaré--Bendixson (\cref{thm:oc-pb}).}

Una órbita periódica es hiperbólica si su multiplicador (la derivada de su mapa de Poincaré, \cref{sec:bf-poincare}) no tiene valor absoluto~$1$. El teorema dice dónde mirar: un retrato del plano solo puede cambiar de forma cuando un equilibrio o una órbita periódica \emph{pierden la hiperbolicidad}, o cuando aparece una conexión entre sillas. Las dos primeras situaciones son locales y ocupan casi todo el capítulo; la tercera es global y se menciona en la \cref{sec:bf-ciclos}.

\begin{definicion}[label=def:bf-bifurcacion]{Bifurcación}
El valor $\mu_c$ es un \clave{valor de bifurcación} de \eqref{eq:bf-familia} si $f(\cdot,\mu_c)$ no es estructuralmente estable: hay valores de $\mu$ arbitrariamente cercanos a $\mu_c$ cuyos retratos no son topológicamente equivalentes al de $\mu_c$. El cambio cualitativo que ocurre al cruzar $\mu_c$ es una \clave{bifurcación}. La gráfica de los equilibrios (o de la amplitud de las órbitas periódicas) en función de $\mu$, con su estabilidad, es el \clave{diagrama de bifurcación}.
\end{definicion}

La \cref{fig:bf-mapa} muestra el recorrido del capítulo. Empezamos en la recta, donde todo se ve (\S\ref{sec:bf-sillanodo}--\ref{sec:bf-transcritica}); pasamos al plano usando la variedad central del \cref{cap:bidimensionales}, que muestra por qué el catálogo de la recta sirve también allí (\S\ref{sec:bf-plano}); vemos la única bifurcación genuinamente bidimensional de un equilibrio, la de Hopf (\S\ref{sec:bf-hopf}), y terminamos con las bifurcaciones de las órbitas periódicas y la herramienta para estudiarlas: el mapa de Poincaré (\S\ref{sec:bf-ciclos}--\ref{sec:bf-poincare}).

\begin{figure}[ht]
\centering
\includegraphics[width=0.95\linewidth]{cap06/t_mapa_bif}
\caption{Mapa del capítulo. Una bifurcación local es una pérdida de hiperbolicidad; en $\R^2$ hay dos maneras de perderla, y la no linealidad decide qué pasa después.}
\label{fig:bf-mapa}
\end{figure}

\notamargen{\textbf{Ruta de las clases.} Semana~2 (10 de agosto): bifurcaciones en la recta, \S\ref{sec:bf-hiperbolicidad}--\ref{sec:bf-transcritica}. Semana~6 (9 de septiembre): introducción a los mapas de Poincaré. Semana~7 (14 de septiembre): bifurcaciones en el plano y Hopf, \S\ref{sec:bf-plano}--\ref{sec:bf-ciclos}; (16 de septiembre): mapas de Poincaré, \S\ref{sec:bf-poincare}.}

\begin{historia}[Una palabra de Poincaré, un teorema de otro Hopf]
El término \emph{bifurcación} lo introdujo Henri Poincaré en 1885, en su memoria sobre las figuras de equilibrio de una masa fluida en rotación~\citep{Poi1885}: al aumentar la velocidad de giro, la familia de elipsoides de equilibrio se \emph{divide} en dos. El pandeo de la columna es aún más viejo: Euler calculó la carga crítica en 1744. La bifurcación que lleva el nombre de Hopf la estudió Poincaré para sistemas planos, Andronov la usó en los años treinta para explicar cómo nacen las oscilaciones en los circuitos de radio~\citep{And29}, y Eberhard Hopf la probó en 1942 en dimensión arbitraria~\citep{Hop42}. Por eso Arnold y la escuela rusa la llaman de Andronov--Hopf. Cuidado con el nombre: el Hopf del teorema del índice del \cref{cap:orbitas} es Heinz Hopf~\citep{Hop35}, un topólogo; el de las bifurcaciones es Eberhard, un analista. No eran parientes.
\end{historia}

\begin{pausa}
Antes de seguir, apueste: un sistema lineal $\dot x=A(\mu)x$, cuyos coeficientes dependen de $\mu$, ¿puede tener bifurcaciones? ¿Puede nacer en él un ciclo límite?
\respuesta{Sí puede cambiar de forma: si un valor propio de $A(\mu)$ cruza el eje imaginario, el origen pasa, por ejemplo, de nodo estable a silla, o de foco estable a foco inestable, y en el cruce aparece una recta de equilibrios o un centro. Pero no puede nacer un ciclo límite: en un sistema lineal, si hay una órbita cerrada hay una familia continua de ellas (\cref{prop:oc-lineal}), y ninguna es aislada. Las bifurcaciones lineales son, además, degeneradas: en el valor crítico no queda claro qué pasa, y son los términos no lineales los que deciden. Eso es este capítulo.}
\end{pausa}

%% ---------------------------------------------------------------------
\section{Qué es exactamente lo que se rompe}
\label{sec:bf-hiperbolicidad}
%% ---------------------------------------------------------------------

Sea $\eq_0$ un equilibrio de \eqref{eq:bf-familia} para $\mu=\mu_0$. ¿Sigue habiendo un equilibrio cerca de $\eq_0$ cuando $\mu$ se mueve un poco? La respuesta la da el teorema de la función implícita aplicado a $f(x,\mu)=0$.

\begin{proposicion}[label=prop:bf-implicita]{Los equilibrios hiperbólicos se continúan}
Si $f$ es de clase $\Cont1$ y $\jacat fx{(\eq_0,\mu_0)}$ es invertible (ningún valor propio nulo), existe una única función suave $\eq(\mu)$, definida para $\mu$ cerca de $\mu_0$, con $\eq(\mu_0)=\eq_0$ y $f(\eq(\mu),\mu)=0$. Si además $\eq_0$ es hiperbólico, $\eq(\mu)$ lo es para $\mu$ cerca de $\mu_0$ y conserva su tipo (mismo número de valores propios con parte real positiva y negativa).
\end{proposicion}

\begin{demostracion}
La primera parte es el teorema de la función implícita. Para la segunda: los valores propios de $\jacat fx{(\eq(\mu),\mu)}$ dependen continuamente de $\mu$, así que ninguno puede pasar de un semiplano al otro sin tocar el eje imaginario, lo que no ocurre cerca de $\mu_0$ porque en $\mu_0$ están a una distancia positiva de él. Con Hartman--Grobman (\cref{thm:bi-hg}), el retrato local no cambia.
\end{demostracion}

La consecuencia es que, cerca de un equilibrio, el retrato solo puede cambiar de forma en un valor $\mu_c$ en el que el equilibrio \alerta{deja de ser hiperbólico}: algún valor propio de la jacobiana toca el eje imaginario. Eso es, por definición, una \clave{bifurcación local}. En $\R^2$, con $\tau=\tr J$ y $\delta=\det J$ (\cref{thm:bi-trazadet}), solo hay dos maneras de tocar el eje imaginario:
\begin{description}
  \item[Ruta 1] Un valor propio real pasa por cero: $\delta(\mu)\to0$. Por aquí ocurren las bifurcaciones silla-nodo, transcrítica y de horquilla.
  \item[Ruta 2] Un par complejo conjugado cruza el eje: $\tau(\mu)\to0$ con $\delta>0$. Por aquí ocurre, y solo ella, la bifurcación de Hopf.
\end{description}
Buscar bifurcaciones locales en el plano es, entonces, seguir el punto $(\tau(\mu),\delta(\mu))$ de cada equilibrio en el plano traza--determinante y ver cuándo cruza el eje $\delta=0$ o la semirrecta $\{\tau=0,\ \delta>0\}$ (\cref{fig:bf-rutas}). No hace falta reconocer formas normales de memoria.

\begin{figure}[ht]
\centering
\includegraphics[width=0.85\linewidth]{cap06/f_rutas}
\caption{El plano traza--determinante del \cref{cap:bidimensionales}, leído como mapa de bifurcaciones. La Ruta~1 cruza el eje $\delta=0$ (un valor propio real pasa por cero); la Ruta~2 cruza la semirrecta $\tau=0$, $\delta>0$ (un par complejo cruza el eje imaginario). Los caminos son los de la silla-nodo de la \cref{sec:bf-plano} y del circuito de la \cref{sec:bf-hopf}.}
\label{fig:bf-rutas}
\end{figure}

\begin{receta}[Cinco pasos para encontrar y clasificar una bifurcación]
\begin{enumerate}
  \item \textbf{Equilibrios en función de $\mu$.} Resolver $f(\eq,\mu)=0$; contar cuántos hay y para qué $\mu$ aparecen o desaparecen.
  \item \textbf{Jacobiano en cada rama.} $J(\mu)=\jacat fx{(\eq(\mu),\mu)}$ y sus invariantes $\tau(\mu)$, $\delta(\mu)$.
  \item \textbf{Cruce.} Resolver $\delta(\mu)=0$ (Ruta~1) o $\tau(\mu)=0$ con $\delta>0$ (Ruta~2). Eso da $\mu_c$.
  \item \textbf{Transversalidad.} Verificar que el cruce es limpio: el valor propio que cruza lo hace con velocidad no nula, $\dif\lambda/\dif\mu\neq0$ (Ruta~1, en la forma de la \cref{sec:bf-sillanodo}) o $\dif(\Re\lambda)/\dif\mu\neq0$ (Ruta~2).
  \item \textbf{No degeneración.} El primer término no lineal que no se anula decide el tipo: cuadrático, silla-nodo; cúbico con simetría, horquilla; en Hopf, el signo del primer coeficiente de Lyapunov.
\end{enumerate}
Con los pasos 1 a 3 se sabe que ``hay algo en $\mu_c$''. Sin el 4 y el 5 no se sabe \emph{qué} hay, ni si el sistema queda a salvo o se va lejos. Casi todos los errores de este tema consisten en nombrar la bifurcación por el parecido del dibujo y saltarse esos dos pasos.
\end{receta}

\begin{trampa}[label=tr:bf-det]{``$\det J=0$ en $\mu_c$: hay una bifurcación''}
\plausible{La \cref{prop:bf-implicita} dice que sin valor propio nulo no hay bifurcación; entonces con valor propio nulo sí la hay.}
La proposición da una condición \emph{necesaria}, no suficiente. En $\dot x=\mu-x^3$, en $(x,\mu)=(0,0)$ la derivada $f_x=-3x^2$ se anula. Sin embargo, para todo $\mu$ hay un único equilibrio, $\eq=\mu^{1/3}$, que es siempre asintóticamente estable (para $\mu\neq0$, $f_x(\eq)=-3\mu^{2/3}<0$; para $\mu=0$, $\dot x=-x^3$). Los retratos para $\mu<0$, $\mu=0$ y $\mu>0$ son todos ``un punto que atrae'': topológicamente iguales. Falló el paso 4: $f_\mu=1\neq0$ pero el término $f_{xx}$ es cero, y el $f_{xxx}$ no basta para crear o destruir equilibrios sin simetría.
\correcto{$\delta(\mu_c)=0$ (o $\tau(\mu_c)=0$ con $\delta>0$) es necesario para una bifurcación local de un equilibrio. Para afirmar que la hay, y de qué tipo, se verifican la transversalidad y la no degeneración.}
\end{trampa}

\begin{pausa}
En $\R^3$, ¿cuántas ``rutas'' hay para perder la hiperbolicidad con un solo parámetro?
\respuesta{Las mismas dos: con un parámetro, lo genérico es que un valor propio real pase por cero o que un par complejo cruce el eje imaginario; los demás valores propios, en general, se quedan lejos. Que pasen dos cosas a la vez (por ejemplo, un cero doble) requiere ajustar dos parámetros: son bifurcaciones de codimensión dos (Bogdanov--Takens, cúspide), fuera del alcance de este curso \citep{Kuz04}. Lo que sí cambia en $\R^3$ es lo que puede pasar con las órbitas periódicas: además de las bifurcaciones de la \cref{sec:bf-ciclos}, aparece la duplicación de período, una de las rutas al caos.}
\end{pausa}

%% ---------------------------------------------------------------------
\section{La bifurcación silla-nodo}
\label{sec:bf-sillanodo}
%% ---------------------------------------------------------------------

Empezamos en la recta, $\dot x=f(x,\mu)$ con $x\in\R$, donde la jacobiana es el número $f_x=\partial f/\partial x$ y la única ruta es que $f_x$ pase por cero. La \clave{bifurcación silla-nodo} es el mecanismo básico por el que se crean o se destruyen equilibrios: al variar el parámetro, dos equilibrios se acercan, chocan y desaparecen.

\begin{ejemplo}[label=ej:bf-sn]{La forma normal $\dot x=\mu+x^2$}
Describa el flujo de $\dot x=\mu+x^2$ para $\mu<0$, $\mu=0$ y $\mu>0$, y dibuje el diagrama de bifurcación.
\solucion La gráfica de $f$ es una parábola que sube con $\mu$ (\cref{fig:bf-sn}, izquierda). Para $\mu<0$ corta el eje en $\eq=\pm\sqrt{-\mu}$; como $f_x(\eq)=2\eq$, el equilibrio $-\sqrt{-\mu}$ es estable y $+\sqrt{-\mu}$ inestable. En $\mu=0$ los dos se funden en $\eq=0$, que es \emph{semiestable}: $\dot x=x^2\ge0$ atrae desde la izquierda y repele hacia la derecha. Para $\mu>0$, $\dot x>0$ en toda la recta: no hay equilibrios y toda solución crece (y escapa en tiempo finito, como en el \cref{ej:un-escape}). En el diagrama (\cref{fig:bf-sn}, derecha) los equilibrios forman la parábola $\mu=-x^2$, con la rama inferior estable (línea continua) y la superior inestable (discontinua).
\end{ejemplo}

\begin{figure}[ht]
\centering
\includegraphics[width=\linewidth]{cap06/f_sillanodo}
\caption{$\dot x=\mu+x^2$. Izquierda: el campo para tres valores de $\mu$, con la línea de fase debajo (círculo lleno: estable; vacío: inestable; medio lleno: semiestable). Derecha: el diagrama de bifurcación, con el parámetro en el eje horizontal y el estado en el vertical, que es la convención.}
\label{fig:bf-sn}
\end{figure}

\codigorepo{cap06_bifurcaciones/C6_Ej1_SillaNodo.m}{sillanodo}{Líneas de fase y diagrama de bifurcación de la silla-nodo}

Con $\dot x=\mu-x^2$ ocurre lo mismo en el otro sentido: al \emph{aumentar} $\mu$ por encima de cero aparecen de la nada dos equilibrios, $\pm\sqrt\mu$. Abraham y Shaw la llamaron ``bifurcación de cielo azul'' \citep[\S3.1]{Str15}, porque los equilibrios aparecen ``out of the clear blue sky''. Es la misma silla-nodo recorrida al revés; no la confunda con la catástrofe de cielo azul de la \cref{sec:bf-ciclos}, que es otra cosa.

\subsection*{¿Por qué siempre una parábola?}

Toda silla-nodo, vista de cerca, se parece a la forma normal. Supongamos que en $(\eq_0,\mu_0)$ se cumple $f=0$ y $f_x=0$. Expandiendo en serie de Taylor alrededor de ese punto,
\[
  \dot x=\underbrace{f}_{=0}+\underbrace{f_x}_{=0}\,(x-\eq_0)+f_\mu\,(\mu-\mu_0)+\tfrac12f_{xx}\,(x-\eq_0)^2+\cdots ,
\]
y si los dos coeficientes que quedan no se anulan, los términos que siguen son pequeños frente a ellos.

\begin{teorema}[label=thm:bf-sn]{Silla-nodo en la recta \citep[\S3.1]{Str15}, \citep[\S3.2]{Kuz04}}
Sea $f$ suave con $f(\eq_0,\mu_0)=0$ y $f_x(\eq_0,\mu_0)=0$. Si además
\[
  a=f_\mu(\eq_0,\mu_0)\neq0\quad\text{(transversalidad)},\qquad b=\tfrac12f_{xx}(\eq_0,\mu_0)\neq0\quad\text{(no degeneración)},
\]
entonces cerca de $(\eq_0,\mu_0)$ el sistema es topológicamente equivalente a $\dot y=a(\mu-\mu_0)+by^2$: hay dos equilibrios hiperbólicos de un lado de $\mu_0$ (donde $a(\mu-\mu_0)/b<0$), uno estable y otro inestable, que se funden en $\mu_0$ y desaparecen del otro lado.
\end{teorema}

\begin{ejemplo}[label=ej:bf-cosecha]{La pesca que colapsa}
En el \cref{ej:un-cosecha} la población con cosecha constante era $\dot x=x(1-x)-h$. Verifique las condiciones del \cref{thm:bf-sn} y escriba la forma normal.
\solucion $f=x-x^2-h$; $f_x=1-2x=0$ en $\eq_0=\tfrac12$, y $f(\tfrac12,h)=0$ da $h_0=\tfrac14$. Además $f_h=-1\neq0$ y $\tfrac12f_{xx}=-1\neq0$. Aquí la expansión es exacta, porque $f$ es cuadrática:
\[
  \dot x=\Bigl(\tfrac14-h\Bigr)-\Bigl(x-\tfrac12\Bigr)^2 .
\]
Para $h<\tfrac14$ hay dos equilibrios, $\tfrac12\pm\sqrt{\tfrac14-h}$ (el mayor estable); en $h=\tfrac14$ chocan, y para $h>\tfrac14$ la población se extingue desde cualquier condición inicial.
\end{ejemplo}

El diodo túnel del \cref{cap:unidimensionales} tiene dos sillas-nodo, una en cada extremo del intervalo de $E$ con tres equilibrios, y entre las dos producen la histéresis de la \cref{fig:un-histeresis}. Los dos fenómenos que siguen, el fantasma y la lentitud crítica, también están en ese circuito.

\subsection*{El fantasma de un equilibrio}

Del lado de la silla-nodo en que no hay equilibrios, el campo sigue siendo pequeño donde estaban: la parábola casi toca el eje. Las soluciones pasan por allí, pero despacio, y tardan mucho en cruzar ese \clave{cuello de botella}. Para $\dot x=-(m+x^2)$, con $m=-\mu>0$, el tiempo para atravesar toda la recta es
\begin{equation}
  T(m)=\int_{-\infty}^{\infty}\frac{\dif x}{m+x^2}=\frac{\pi}{\sqrt m} ,
  \label{eq:bf-fantasma}
\end{equation}
y casi todo ese tiempo se gasta cerca de $x=0$, donde estaba el equilibrio (\cref{fig:bf-fantasma}). Con $m=0.01$ la solución tarda $2\arctan(10)/0.1=29.42$ unidades de tiempo solo en ir de $x=1$ a $x=-1$. Es el \clave{fantasma} de la silla-nodo: no hay equilibrio, pero el sistema se comporta durante un rato largo como si lo hubiera.

\begin{figure}[ht]
\centering
\includegraphics[width=\linewidth]{cap06/f_fantasma}
\caption{El fantasma. Izquierda: $\dot x=\mu-x^2$ desde $x(0)=4$ para cuatro valores de $\mu<0$; cuanto más cerca de $\mu=0$, más larga la pausa cerca de $x=0$. Derecha: el tiempo de tránsito de $x=1$ a $x=-1$, simulado (puntos) y la ley $\pi/\sqrt{\abs\mu}$ (recta de pendiente $-\tfrac12$ en escala log--log).}
\label{fig:bf-fantasma}
\end{figure}

\codigorepo{cap06_bifurcaciones/C6_Ej2_Fantasma.m}{fantasma}{El tiempo de tránsito por el cuello de botella contra la ley $\pi/\sqrt{\abs\mu}$}

\video{V6_1_cuello_botella}{El fantasma de la silla-nodo}

\begin{trampa}[label=tr:bf-fantasma]{``Ya no hay equilibrio: el estado pasa de largo''}
\plausible{Para $\mu$ apenas al otro lado del valor crítico no hay equilibrio, así que la solución no tiene dónde detenerse y atraviesa la región sin demora.}
La solución no se detiene, pero tarda un tiempo que crece como $1/\sqrt{\abs{\mu-\mu_c}}$, sin cota. En una simulación corta, o en una medición, el sistema parece estar en reposo, y se puede confundir el cuello de botella con un equilibrio. Lo mismo pasa en el oscilador no uniforme del \cref{cap:orbitas}: su período \eqref{eq:oc-periodo} crece como $1/\sqrt{\omega-a}$ porque la fase pasa por el fantasma de una silla-nodo en el círculo.
\correcto{La divergencia $1/\sqrt{\abs{\mu-\mu_c}}$ del tiempo de tránsito es la firma de una silla-nodo cercana, y es una predicción cuantitativa que se puede verificar con una simulación (\cref{fig:bf-fantasma}).}
\end{trampa}

Del otro lado ocurre algo parecido. Para $\dot x=\mu-x^2$ con $\mu>0$, el equilibrio estable $\sqrt\mu$ tiene $f_x=-2\sqrt\mu$: la constante de tiempo de la recuperación ante una perturbación es $1/(2\sqrt\mu)$, que vale $0.5$ con $\mu=1$ y $5$ con $\mu=0.01$. Al acercarse a la silla-nodo el sistema se vuelve cada vez más lento para recuperarse. Esa \clave{lentitud crítica} se usa como señal de alerta temprana de transiciones abruptas en ecología, clima y redes eléctricas: la autocorrelación y la varianza de las fluctuaciones crecen antes de que el equilibrio desaparezca~\citep{SBBB09}.

\begin{pausa}
Para $\dot x=\mu-x^2$ con $\mu<0$, ¿es $\pi/\sqrt{\abs\mu}$ el tiempo que tarda una solución en ir de $x=+\infty$ a $x=-\infty$? ¿Tiene sentido que sea finito?
\respuesta{Sí: \eqref{eq:bf-fantasma} es la integral sobre toda la recta, y es finita porque lejos del origen $\abs{\dot x}\approx x^2$, que crece tan rápido que la solución ``llega del infinito'' y ``se va al infinito'' en tiempo finito (es el escape en tiempo finito del \cref{cap:unidimensionales}, hacia atrás y hacia adelante). Casi todo el tiempo se gasta cerca del fantasma: con $\abs\mu=0.01$, de los $31.42$ que dura todo el tránsito, $29.42$ se gastan entre $x=1$ y $x=-1$.}
\end{pausa}

%% ---------------------------------------------------------------------
\section{Transcrítica y horquilla, y por qué son frágiles}
\label{sec:bf-transcritica}
%% ---------------------------------------------------------------------

\subsection{Intercambio de estabilidad: la transcrítica}

Hay modelos en los que un equilibrio existe para todo valor del parámetro: el número de fotones en un láser, una población o una concentración nula siguen siéndolo pase lo que pase. Ese equilibrio no puede desaparecer en una silla-nodo, pero sí puede cambiar de estabilidad. La forma normal es
\begin{equation}
  \dot x=\mu x-x^2 .
  \label{eq:bf-tc}
\end{equation}
Sus equilibrios son $\eq=0$ y $\eq=\mu$, con $f_x(0)=\mu$ y $f_x(\mu)=-\mu$. Para $\mu<0$ el origen es estable y $\mu$ inestable; para $\mu>0$ es al revés. Las dos ramas se cruzan en $\mu=0$ e \alerta{intercambian su estabilidad} (\cref{fig:bf-tc}): es la \clave{bifurcación transcrítica}. A diferencia de la silla-nodo, ningún equilibrio desaparece.

\begin{ejemplo}[label=ej:bf-laser]{El umbral de un láser}
En un láser de estado sólido, una fuente externa ``bombea'' los átomos a un estado excitado; cada fotón que pasa estimula a un átomo excitado a emitir otro (ganancia), y los fotones escapan por los espejos (pérdida). Con $n(t)$ el número de fotones, $N(t)$ el de átomos excitados, $G>0$ la ganancia y $k>0$ la tasa de pérdida, $\dot n=GnN-kn$. Cada emisión devuelve un átomo a su estado base, así que $N$ baja con $n$: en un modelo simplificado, $N=N_0-\alpha n$, con $N_0$ el número de átomos excitados por el bombeo cuando no hay acción láser \citep[\S3.3]{Str15}. Encuentre el umbral del láser.
\solucion Reemplazando,
\[
  \dot n=(GN_0-k)\,n-\alpha G\,n^2 ,
\]
que es \eqref{eq:bf-tc} con parámetro $GN_0-k$ y un cambio de escala. $n=0$ es equilibrio para todo $N_0$, con $f_n(0)=GN_0-k$. Si $N_0<k/G$, el origen es estable: no hay emisión estimulada y el dispositivo es una lámpara. En $N_0=k/G$ hay una transcrítica, y para $N_0>k/G$ el origen es inestable y aparece el equilibrio estable $n^*=(GN_0-k)/(\alpha G)>0$: el láser emite. $N_0=k/G$ es el \clave{umbral del láser}. El otro equilibrio, $n^*<0$ para $N_0<k/G$, no tiene sentido físico, pero su existencia es la que obliga a la transcrítica.
\end{ejemplo}

\begin{figure}[ht]
\centering
\includegraphics[width=\linewidth]{cap06/f_transcritica}
\caption{Izquierda: $\dot x=\mu x-x^2$, con las dos ramas que se cruzan e intercambian la estabilidad. Derecha: el umbral del láser, $n^*$ contra el bombeo $N_0$ (con $G=k=1$, $\alpha=\tfrac12$); la rama negativa es la que no tiene sentido físico.}
\label{fig:bf-tc}
\end{figure}

\codigorepo{cap06_bifurcaciones/C6_Ej3_Transcritica.m}{laser}{El láser: número de fotones en régimen permanente contra el bombeo}

\subsection{Simetría: las horquillas}

Cuando el problema tiene una simetría izquierda--derecha, como la columna que puede pandearse hacia cualquiera de los dos lados, los equilibrios aparecen y desaparecen en pares simétricos. La forma normal de la \clave{horquilla supercrítica} es
\begin{equation}
  \dot x=\mu x-x^3 ,
  \label{eq:bf-pf}
\end{equation}
que no cambia al reemplazar $x$ por $-x$: esa invariancia es la expresión matemática de la simetría. Para $\mu<0$ el origen es el único equilibrio y es estable; para $\mu>0$ es inestable y aparecen dos equilibrios estables, $\pm\sqrt\mu$, con $f_x(\pm\sqrt\mu)=-2\mu$ (\cref{fig:bf-pf}, izquierda). En $\mu=0$ el origen es estable pero no exponencialmente: $\dot x=-x^3$ decae como $1/\sqrt t$.

Si el término cúbico desestabiliza, $\dot x=\mu x+x^3$, se tiene la \clave{horquilla subcrítica}: los equilibrios $\pm\sqrt{-\mu}$ existen para $\mu<0$ y son \emph{inestables} ($f_x=-2\mu>0$); son el borde de la cuenca del origen, $(-\sqrt{-\mu},\sqrt{-\mu})$, que se encoge y desaparece en $\mu=0$. Para $\mu>0$ todas las soluciones con $x_0\neq0$ escapan, en tiempo finito en esta forma normal truncada. En un sistema físico algo detiene el crecimiento; el modelo más simple agrega un término quíntico estabilizante:
\begin{equation}
  \dot x=\mu x+x^3-x^5 .
  \label{eq:bf-quintica}
\end{equation}
Sus equilibrios no nulos cumplen $x^2=\tfrac12\bigl(1\pm\sqrt{1+4\mu}\bigr)$, y hay dos sillas-nodo simétricas en $\mu=-\tfrac14$, $x=\pm1/\sqrt2$ (\cref{fig:bf-pf}, derecha). Para $-\tfrac14<\mu<0$ coexisten tres estados estables: el origen y las dos ramas grandes. Si $\mu$ sube desde valores negativos, el estado se queda en el origen hasta $\mu=0$ y allí \alerta{salta} a una rama grande; si después $\mu$ baja, se queda en la rama grande hasta $\mu=-\tfrac14$ y solo entonces cae al origen. Es una \clave{histéresis}, la misma del diodo túnel. Con $\mu=-0.1$, la rama estable está en $x=0.9420$ y la inestable en $x=0.3357$.

\begin{figure}[ht]
\centering
\includegraphics[width=\linewidth]{cap06/f_horquillas}
\caption{Izquierda: horquilla supercrítica $\dot x=\mu x-x^3$. Centro: subcrítica $\dot x=\mu x+x^3$; las ramas inestables limitan la cuenca del origen. Derecha: subcrítica con término quíntico \eqref{eq:bf-quintica}; las flechas recorren el ciclo de histéresis entre $\mu=-\tfrac14$ y $\mu=0$.}
\label{fig:bf-pf}
\end{figure}

\codigorepo{cap06_bifurcaciones/C6_Ej4_Horquillas.m}{quintica}{Barrido lento de $\mu$ en la horquilla subcrítica con término quíntico: el ciclo de histéresis}

\begin{trampa}[label=tr:bf-subcritica]{``Pasado $\mu=0$, el estado se va a la rama $\pm\sqrt{-\mu}$''}
\plausible{En la supercrítica, al cruzar $\mu=0$ el estado pasa del origen a una de las ramas que nacen; en la subcrítica debe pasar algo parecido.}
En la subcrítica las ramas $\pm\sqrt{-\mu}$ existen del \emph{otro} lado ($\mu<0$), son inestables y se funden con el origen justo en $\mu=0$. Pasado ese valor no queda ningún atractor cerca: el estado se va lejos, al infinito en la forma normal o a otra rama (la quíntica) en un sistema con más términos. Y no regresa al bajar $\mu$ un poco.
\correcto{Supercrítica: el atractor cambia continuamente con $\mu$ (bifurcación \clave{segura} o \emph{suave}). Subcrítica: el atractor desaparece y el estado salta a otro lejano (bifurcación \clave{peligrosa} o \emph{dura}), casi siempre con histéresis \citep[\S2.7]{Kha02}.}
\end{trampa}

\subsection{Cuáles sobreviven a una perturbación}

Los modelos nunca son exactos. ¿Qué pasa si al campo se le suma una constante pequeña $\varepsilon$, que representa una imperfección? La \cref{fig:bf-imperfectas} lo muestra:
\begin{itemize}
  \item \emph{Silla-nodo:} $\dot x=\mu-x^2+\varepsilon$ es la misma bifurcación, corrida a $\mu=-\varepsilon$. Sobrevive.
  \item \emph{Transcrítica:} $\dot x=\mu x-x^2+\varepsilon$ tiene equilibrios $x=\tfrac12\bigl(\mu\pm\sqrt{\mu^2+4\varepsilon}\bigr)$. Si $\varepsilon>0$ el discriminante es siempre positivo y las dos ramas ya no se tocan; si $\varepsilon<0$ se convierte en dos sillas-nodo, en $\mu=\pm2\sqrt{-\varepsilon}$ ($\pm0.69$ con $\varepsilon=-0.12$). Se deshace.
  \item \emph{Horquilla:} $\dot x=\mu x-x^3+h$ se parte en una rama continua y una silla-nodo aislada, en $\mu=3(h/2)^{2/3}$ ($0.4072$ con $h=0.1$). Se deshace.
\end{itemize}
Solo la silla-nodo (y, en el plano, la de Hopf) es \clave{genérica}: aparece en una familia de un parámetro cualquiera y sobrevive a perturbaciones. La transcrítica y la horquilla necesitan una estructura del modelo que \emph{prohíba} la perturbación que las destruiría: la transcrítica, que $x=0$ sea un equilibrio para todo $\mu$ (por eso aparece en poblaciones, concentraciones y fotones, donde el cero es físico); la horquilla, la simetría $x\mapsto-x$. Una columna real nunca es perfectamente recta, y por eso se pandea siempre hacia el mismo lado, el de su imperfección.

\begin{figure}[ht]
\centering
\includegraphics[width=\linewidth]{cap06/f_imperfectas}
\caption{Bifurcaciones imperfectas. Gris: sin perturbación. La silla-nodo solo se corre; la transcrítica se separa en dos ramas que no se tocan ($\varepsilon>0$) o en dos sillas-nodo ($\varepsilon<0$); la horquilla se parte en una rama continua y una silla-nodo aislada.}
\label{fig:bf-imperfectas}
\end{figure}

\codigorepo{cap06_bifurcaciones/C6_Ej5_Imperfectas.m}{imperfecta}{Las tres bifurcaciones con una imperfección $\varepsilon$}

\begin{trampa}[label=tr:bf-generica]{``La transcrítica y la horquilla son tan comunes como la silla-nodo''}
\plausible{Las tres aparecen en el mismo catálogo, con formas normales igual de simples.}
Las formas normales son simples porque ya traen la estructura: $\dot x=\mu x-x^2$ deja fijo el cero y $\dot x=\mu x-x^3$ es impar. Sin esa estructura, una perturbación arbitraria las convierte en sillas-nodo o en ramas que no se cruzan.
\correcto{Antes de nombrar una transcrítica o una horquilla, identifique qué propiedad del modelo la protege: un equilibrio que existe para todo $\mu$, o una simetría. Si no la encuentra, lo que tiene es (casi seguro) una silla-nodo o una imperfección.}
\end{trampa}

\begin{pausa}
En el modelo del láser, ¿por qué la bifurcación es transcrítica y no silla-nodo? ¿Qué pasaría si hubiera un poco de emisión espontánea, $\dot n=(GN_0-k)n-\alpha Gn^2+\varepsilon$ con $\varepsilon>0$ pequeño?
\respuesta{Porque $n=0$ es equilibrio para todo $N_0$: sin fotones no hay emisión estimulada. Con emisión espontánea, $n=0$ deja de ser equilibrio ($\dot n=\varepsilon>0$), la transcrítica se deshace en dos ramas separadas, y el umbral deja de ser nítido: la potencia de salida crece de manera suave, aunque rápida, alrededor de $N_0=k/G$. Es lo que se mide en un láser real.}
\end{pausa}

%% ---------------------------------------------------------------------
\section{Bifurcaciones de equilibrios en el plano}
\label{sec:bf-plano}
%% ---------------------------------------------------------------------

\subsection{Por qué el catálogo de la recta sirve en el plano}

En la Ruta~1 un valor propio pasa por cero y el otro, digamos, se queda negativo. Es exactamente la situación de la variedad central (\cref{thm:bi-central}): la dirección estable se ``apaga'' exponencialmente rápido y la dinámica que decide ocurre sobre una curva invariante $W^c$, donde es escalar. Para incluir el parámetro se usa un truco: se agrega la ecuación $\dot\mu=0$. El sistema aumentado $(x,\mu)$ tiene un valor propio nulo más (el de $\mu$), y su variedad central tiene dimensión dos: es una superficie $x_s=h(x_c,\mu)$ sobre la que vive una ecuación escalar con parámetro, $\dot x_c=g(x_c,\mu)$ \citep[\S3.2]{GuHo83}. A esa ecuación se le aplica el catálogo de las secciones anteriores.

\begin{ideaclave}
Una bifurcación de la Ruta~1 en el plano es una bifurcación de la recta que ocurre sobre la variedad central, mientras la otra dirección contrae. Por eso las formas normales en el plano se escriben desacopladas, con $\dot x_2=-x_2$: allí $W^c$ es exactamente el eje $x_1$. En un sistema real $W^c$ es una curva que hay que calcular; el catálogo sigue valiendo, pero los coeficientes cambian.
\end{ideaclave}

\subsection{La silla-nodo en el plano}

La forma normal es
\begin{equation}
  \dot x_1=\mu-x_1^2,\qquad \dot x_2=-x_2 ,\qquad J=\begin{bmatrix}-2x_1&0\\0&-1\end{bmatrix} .
  \label{eq:bf-sn2}
\end{equation}
Para $\mu>0$ hay dos equilibrios, $(\pm\sqrt\mu,0)$. En $(\sqrt\mu,0)$, $\tau=-2\sqrt\mu-1$ y $\delta=2\sqrt\mu>0$: nodo estable. En $(-\sqrt\mu,0)$, $\delta=-2\sqrt\mu<0$: silla. Al bajar $\mu$, el nodo y la silla se acercan sobre el eje $x_1$, chocan en $\mu=0$ y desaparecen; de ahí el nombre ``silla-nodo'' (\cref{fig:bf-sn2}). En el plano traza--determinante, los dos puntos $(\tau,\delta)$ llegan a $(-1,0)$ desde lados opuestos del eje $\delta=0$: es el cruce de la Ruta~1 (\cref{fig:bf-rutas}). El valor propio que va a cero, $\lambda_1=-2\sqrt\mu$, es el de la dirección $x_1$; el otro se queda en $-1$.

\begin{figure}[ht]
\centering
\includegraphics[width=\linewidth]{cap06/f_sillanodo2d}
\caption{Silla-nodo en el plano, \eqref{eq:bf-sn2}, para $\mu=-1$, $0$ y $1$, con las nulclinas de $x_1$ (verde) y de $x_2$ (ámbar). Para $\mu<0$ las trayectorias se frenan cerca del origen, donde está el fantasma, y siguen de largo.}
\label{fig:bf-sn2}
\end{figure}

\codigorepo{cap06_bifurcaciones/C6_Ej6_SillaNodo2D.m}{sillanodo2d}{Retratos de fase y el recorrido de $(\tau,\delta)$ de cada rama}

\begin{trampa}[label=tr:bf-traza]{``En $\mu=0$, $\tau=-1<0$: el origen es estable''}
\plausible{La traza es negativa y el determinante no es negativo; por la clasificación del \cref{cap:bidimensionales}, el origen es estable.}
En $\mu=0$ el determinante es \emph{cero}: hay un valor propio nulo, el equilibrio no es hiperbólico y la linealización no decide (\cref{thm:bi-hg}). El retrato linealizado, $\dot x_1=0$, $\dot x_2=-x_2$, es una recta de equilibrios, cualitativamente falso. Sobre la variedad central, $\dot x_1=-x_1^2$: desde cualquier $x_1(0)<0$ la solución escapa a $-\infty$. El origen es inestable.
\correcto{Con $\delta=0$ la traza no dice nada sobre la estabilidad. Se reduce a la variedad central y se mira el primer término no lineal.}
\end{trampa}

La transcrítica y las horquillas en el plano se escriben igual: $\dot x_1=\mu x_1-x_1^2$ (o $\mu x_1\mp x_1^3$), $\dot x_2=-x_2$. En la transcrítica, el origen pasa de nodo estable a silla y $(\mu,0)$ de silla a nodo estable. En la horquilla supercrítica, para $\mu>0$ el origen es una silla y $(\pm\sqrt\mu,0)$ dos nodos estables. En la subcrítica, $\dot x_1=\mu x_1+x_1^3$, para $\mu<0$ hay tres equilibrios, $(0,0)$ y $(\pm\sqrt{-\mu},0)$: el origen es un nodo estable y los otros dos son sillas, cuyas variedades estables limitan la cuenca del origen; al subir $\mu$ las dos sillas se funden con el origen y para $\mu>0$ solo queda el origen, convertido en silla.

\begin{ejemplo}[label=ej:bf-escondida]{Una horquilla escondida}
Estudie el origen de $\dot x_1=\mu x_1-x_1x_2$, $\dot x_2=-x_2+x_1^2$ para $\mu$ cerca de cero.
\solucion La jacobiana en el origen es $\diag(\mu,-1)$: nodo estable si $\mu<0$, silla si $\mu>0$, y en $\mu=0$ un valor propio nulo. ¿Qué bifurcación es? Si se ``proyecta'' sobre el eje $x_1$ poniendo $x_2=0$, queda $\dot x_1=\mu x_1$, que no decide nada en $\mu=0$. Hay que usar la variedad central. Con $\mu=0$, la ecuación de invariancia \eqref{eq:bi-pde} para $x_2=h(x_1)=c_2x_1^2+c_4x_1^4+\cdots$ es
\[
  h'(x_1)\bigl(-x_1h(x_1)\bigr)=-h(x_1)+x_1^2 ,
\]
y da $c_2=1$, $c_4=2$. Sobre la variedad central, $\dot x_1=\mu x_1-x_1h(x_1)\approx\mu x_1-x_1^3$: una \alerta{horquilla supercrítica}. En este ejemplo se confirma con los equilibrios exactos: $\dot x_1=0$ con $x_1\neq0$ da $x_2=\mu$, y $\dot x_2=0$ da $x_1^2=x_2=\mu$; para $\mu>0$ aparecen $(\pm\sqrt\mu,\mu)$, con $\tau=-1$ y $\delta=2\mu>0$: estables. Nótese que los equilibrios nuevos no están sobre el eje $x_1$ sino sobre la parábola $x_2=x_1^2$, que es la variedad central. La simetría $(x_1,x_2)\mapsto(-x_1,x_2)$ del sistema es la que protege la horquilla.
\end{ejemplo}

\begin{figure}[ht]
\centering
\includegraphics[width=\linewidth]{cap06/f_escondida}
\caption{La horquilla escondida del \cref{ej:bf-escondida} con $\mu=0.1$: el origen es una silla y las trayectorias caen sobre la parábola $x_2\approx x_1^2$ (verde) antes de llegar a $(\pm\sqrt\mu,\mu)$. A la derecha, el diagrama de bifurcación de $x_1^*$.}
\label{fig:bf-escondida}
\end{figure}

\begin{trampa}[label=tr:bf-proyectar]{``Para reducir, se pone $x_2=0$''}
\plausible{$x_2$ es la variable rápida y estable, así que se puede suponer que ya llegó a cero y estudiar $\dot x_1$ sola.}
Es la misma trampa del \cref{cap:bidimensionales} (\cref{tr:bi-zcero}), ahora con un parámetro: con $x_2=0$ se pierde el término $-x_1^3$, que es el que decide el tipo de bifurcación. $x_2$ no tiende a cero; tiende a $h(x_1)\approx x_1^2$. En las formas normales desacopladas proyectar sobre el eje da la respuesta correcta, pero por construcción, no por método.
\correcto{Se reduce a la variedad central, $x_2=h(x_1,\mu)$, calculada con la ecuación de invariancia hasta el orden que decida.}
\end{trampa}

\subsection{Un ejemplo completo: el interruptor genético}

\begin{ejemplo}[label=ej:bf-griffith]{El interruptor genético de Griffith}
En un modelo de control genético, la actividad de un gen es inducida por dos copias de la proteína que el mismo gen produce: es un lazo de realimentación positiva~\citep{Gri68}, \citep[Ejemplo~8.1.2]{Str15}. Con $x$ y $y$ proporcionales a las concentraciones de la proteína y del ARN mensajero,
\begin{equation}
  \dot x=-ax+y,\qquad \dot y=\frac{x^2}{1+x^2}-by ,\qquad a,b>0,
  \label{eq:bf-griffith}
\end{equation}
donde $a$ y $b$ son las tasas de degradación. Muestre que hay tres equilibrios si y solo si $a<a_c$, encuentre $a_c$, muestre que en $a=a_c$ dos de ellos chocan en una silla-nodo e interprete el resultado.
\solucion \emph{Paso 1, equilibrios.} Las nulclinas son la recta $y=ax$ ($\dot x=0$) y la sigmoide $y=x^2/\bigl(b(1+x^2)\bigr)$ ($\dot y=0$). Igualando, $x\bigl[ab(1+x^2)-x\bigr]=0$. Siempre está $\eq=0$; los demás cumplen $abx^2-x+ab=0$:
\[
  \eq_\pm=\frac{1\pm\sqrt{1-4a^2b^2}}{2ab} ,
\]
reales y distintos si y solo si $4a^2b^2<1$, es decir, $a<a_c=\dfrac1{2b}$. En $a=a_c$ coinciden en $\eq=1/(2a_cb)=1$, \emph{sea cual sea} $b$. Geométricamente, la recta es tangente a la sigmoide (\cref{fig:bf-griffith}, izquierda).

\emph{Paso 2, jacobiano.}
\[
  J=\begin{bmatrix}-a&1\\ \dfrac{2x}{(1+x^2)^2}&-b\end{bmatrix},\qquad \tau=-(a+b)<0,\qquad \delta=ab-\frac{2x}{(1+x^2)^2} .
\]
En los equilibrios no nulos, $ab=x/(1+x^2)$, y queda $\delta=\dfrac{x(x^2-1)}{(1+x^2)^2}$.

\emph{Paso 3, el cruce.} $\delta=0$ en $x=1$, que es justo donde chocan los equilibrios en $a=a_c$: Ruta~1 con $\tau<0$. El equilibrio intermedio ($0<\eq_-<1$) tiene $\delta<0$: es una \alerta{silla}. El alto ($\eq_+>1$) tiene $0<\delta<ab$, así que $\tau^2-4\delta>(a+b)^2-4ab=(a-b)^2\ge0$: es un \alerta{nodo estable}. El origen tiene valores propios $-a$ y $-b$: siempre es un nodo estable.

\emph{Pasos 4 y 5.} Al aumentar $a$, el producto $ab$ que fija los equilibrios cambia con velocidad no nula, y las dos raíces de una cuadrática se funden con una tangencia simple (la sigmoide tiene curvatura no nula en $x=1$): es una silla-nodo. Con $b=0.35$, $a_c=1.4286$; con $a=1$ los equilibrios son $0$, $0.4084$ (silla) y $2.4488$ (nodo).

\emph{Interpretación.} Para $a<a_c$ el gen tiene dos estados estables: \emph{apagado} ($\eq=0$) y \emph{encendido} ($\eq_+$), separados por un umbral, que es la variedad estable de la silla. La célula ``recuerda'' en cuál está: es un \clave{interruptor biestable}, un mecanismo de memoria celular. Pero solo existe si la proteína y el mensajero se degradan despacio, $ab<\tfrac12$.
\end{ejemplo}

\begin{figure}[ht]
\centering
\includegraphics[width=\linewidth]{cap06/f_griffith}
\caption{El interruptor genético con $b=0.35$. Izquierda: la sigmoide $y=x^2/\bigl(b(1+x^2)\bigr)$ y la recta $y=ax$ para $a<a_c$ (tres cortes), $a=a_c$ (tangente en $x=1$) y $a>a_c$ (solo el origen). Centro: retrato de fase con $a=1$; la variedad estable de la silla separa las cuencas del gen apagado y encendido. Derecha: diagrama de bifurcación, con el ciclo de histéresis.}
\label{fig:bf-griffith}
\end{figure}

\codigorepo{cap06_bifurcaciones/C6_Ej7_Genetico.m}{genetico}{Equilibrios, tipo y valor crítico del interruptor genético}

\begin{trampa}[label=tr:bf-reversible]{``Si $a$ vuelve a bajar, el gen se vuelve a encender''}
\plausible{Para $a<a_c$ el estado encendido existe; al regresar $a$ a un valor donde existía, la célula debería regresar a él.}
Al subir $a$ por encima de $a_c$ el estado encendido desaparece y la célula cae al apagado. Al bajar $a$ de nuevo, el equilibrio alto reaparece, pero la célula ya está en la cuenca del origen, que es estable para todo $a$, y allí se queda. Con $b=0.35$: con $a=1.3$ el gen encendido sigue encendido; con $a=1.6$ se apaga; de vuelta en $a=1.3$ sigue apagado.
\correcto{Una silla-nodo con otro atractor del otro lado produce histéresis: el estado depende de la historia del parámetro, no solo de su valor actual. Para encender el gen otra vez hay que sacarlo de la cuenca del origen, por ejemplo con un pulso de proteína que cruce la variedad estable de la silla.}
\end{trampa}

\begin{modelo}[InterruptorGenetico.slx]
\cmd{crearModelo\_InterruptorGenetico.m} arma \eqref{eq:bf-griffith} con dos integradores y un bloque \texttt{Fcn} para la sigmoide, y hace variar $a$ lentamente con una rampa que sube de $1$ a $1.8$ y vuelve a bajar. La gráfica de $x$ contra $a$ recorre el ciclo de histéresis de la \cref{fig:bf-griffith}: el gen se apaga en $a\approx a_c$ y no se vuelve a encender.
\end{modelo}

\begin{pausa}
Calcule $a_c$ para el interruptor genético por dos caminos independientes: como el valor donde se anula el discriminante de la cuadrática, y como el valor donde $\delta=0$ sobre la rama de equilibrios. ¿Por qué coinciden?
\respuesta{Por el discriminante: $1-4a^2b^2=0$ da $a_c=1/(2b)$. Por el determinante: $\delta=0$ en $x=1$, y $x=1$ es equilibrio si $ab(1+1)=1$, es decir, $a=1/(2b)$. Coinciden porque las dos condiciones dicen lo mismo: el equilibrio es una raíz doble de $f(x,a)=0$, y una raíz doble es exactamente un punto donde $f=0$ y $\partial f/\partial x=0$ (en el plano, $\det J=0$). Es la \cref{prop:bf-implicita} leída al revés.}
\end{pausa}

%% ---------------------------------------------------------------------
\section{La bifurcación de Hopf}
\label{sec:bf-hopf}
%% ---------------------------------------------------------------------

Hasta aquí un nodo perdía la estabilidad porque un valor propio real pasaba por cero. ¿Qué pasa cuando un \emph{foco} estable pierde la estabilidad, es decir, cuando un par de valores propios complejos conjugados cruza el eje imaginario (Ruta~2)? Ningún equilibrio se crea ni se destruye: $\delta>0$ durante todo el cruce, así que la \cref{prop:bf-implicita} garantiza que el equilibrio sigue ahí, solo y suave en $\mu$. Lo que nace es una \alerta{oscilación}.

\begin{ejemplo}[label=ej:bf-circuito]{Un circuito con un resistor que depende de la temperatura}
Un circuito $LC$ con un resistor no lineal cuya característica $i=f_\mu(v)$ depende de un parámetro $\mu$ (la temperatura, por ejemplo) tiene, en variables adimensionales, la forma de Liénard
\begin{equation}
  \dot x_1=x_2-f_\mu(x_1),\qquad \dot x_2=-x_1,\qquad f_\mu(x_1)=x_1^3-\mu x_1 ,
  \label{eq:bf-circuito}
\end{equation}
con $\mu\in[-1,1]$. Con $\mu=1$ es el oscilador de van der Pol del \cref{cap:orbitas}. Estudie el origen al variar $\mu$.
\solucion El único equilibrio es el origen, con
\[
  J=\begin{bmatrix}\mu&1\\-1&0\end{bmatrix},\qquad \tau=\mu,\quad \delta=1,\quad \lambda_{1,2}=\tfrac12\bigl(\mu\pm\sqrt{\mu^2-4}\bigr) .
\]
Como $\tau^2-4\delta=\mu^2-4<0$ en $[-1,1]$: para $\mu<0$ es un foco estable, para $\mu>0$ un foco inestable, y en $\mu=0$ la linealización es un centro y \emph{no decide}. El par cruza el eje imaginario en $\pm j$ con velocidad $\dif(\Re\lambda)/\dif\mu=\tfrac12\neq0$. Físicamente: para $\mu<0$ el resistor disipa y todo decae al origen; para $\mu>0$ el resistor es activo cerca del origen (resistencia negativa, $f_\mu'(0)=-\mu<0$) y las oscilaciones pequeñas crecen. Pero lejos del origen el término $x_1^3$ vuelve a disipar, y las oscilaciones se estabilizan en un \alerta{ciclo límite} cuya amplitud crece desde cero con $\mu$ (\cref{fig:bf-circuito}).
\end{ejemplo}

\begin{figure}[ht]
\centering
\includegraphics[width=\linewidth]{cap06/f_circuito}
\caption{El circuito \eqref{eq:bf-circuito}. Arriba: retratos para $\mu=-0.5$, $0$, $0.5$ y $1$. Abajo, izquierda: la amplitud del ciclo (máximo de $x_1$) contra $\mu$ en escala log--log, con la predicción $2\sqrt{\mu/3}$ de pendiente $\tfrac12$. Abajo, derecha: el período, que vale $2\pi$ solo en el límite $\mu\to0^+$ y llega a $6.663$ en $\mu=1$.}
\label{fig:bf-circuito}
\end{figure}

\codigorepo{cap06_bifurcaciones/C6_Ej8_HopfCircuito.m}{circuito}{Amplitud y período del ciclo del circuito contra $\mu$}

\subsection{La forma normal}

La forma más simple de ver qué pasa es un sistema que en coordenadas polares se desacopla:
\begin{equation}
  \dot x_1=\mu x_1-x_2-x_1(x_1^2+x_2^2),\quad \dot x_2=x_1+\mu x_2-x_2(x_1^2+x_2^2)
  \quad\Longleftrightarrow\quad
  \dot r=r(\mu-r^2),\quad \dot\theta=1 .
  \label{eq:bf-hopfnormal}
\end{equation}
La jacobiana en el origen es $\bigl[\begin{smallmatrix}\mu&-1\\1&\mu\end{smallmatrix}\bigr]$, con $\lambda=\mu\pm j$. La ecuación radial es una horquilla supercrítica en $r\ge0$: para $\mu\le0$, $\dot r<0$ para todo $r>0$ y el origen es asintóticamente estable; para $\mu>0$ el origen es inestable y $r=\sqrt\mu$ es estable (si $0<r<\sqrt\mu$, $\dot r>0$; si $r>\sqrt\mu$, $\dot r<0$). En el plano: \alerta{una espiral estable se convierte en una espiral inestable rodeada por un ciclo límite estable} de radio $\sqrt\mu$. Es la \clave{bifurcación de Hopf supercrítica}.

Con el signo cambiado, $\dot r=r(\mu+r^2)$, la jacobiana es idéntica, pero para $\mu<0$ el origen estable está rodeado por un ciclo \emph{inestable} de radio $\sqrt{-\mu}$, que es el borde de su cuenca; el ciclo se contrae, se funde con el origen en $\mu=0$, y para $\mu>0$ las soluciones escapan. Es la \clave{Hopf subcrítica} (\cref{fig:bf-supersub}).

\begin{figure}[ht]
\centering
\includegraphics[width=\linewidth]{cap06/f_supersub}
\caption{La misma linealización, destinos opuestos. Izquierda: supercrítica, $\dot r=r(\mu-r^2)$; el ciclo estable nace con amplitud $\sqrt\mu$. Derecha: subcrítica, $\dot r=r(\mu+r^2)$; el ciclo inestable (umbral de seguridad) se contrae hasta el origen y para $\mu>0$ la trayectoria escapa.}
\label{fig:bf-supersub}
\end{figure}

\video{V6_2_hopf}{Nace un ciclo límite}

\subsection{El teorema}

El circuito no está en coordenadas polares, y su ciclo no es un círculo. Lo que dice que la forma normal es representativa es el teorema de Hopf.

\begin{teorema}[label=thm:bf-hopf]{Poincaré--Andronov--Hopf \citep{Hop42}, \citep[\S3.4]{GuHo83}, \citep[\S2.7]{Kha02}}
Sea $\dot x=f(x,\mu)$, $x\in\R^2$, suave, con un equilibrio $\eq(\mu)$ cuya jacobiana tiene valores propios $\alpha(\mu)\pm j\omega(\mu)$. Supongamos que en $\mu=\mu_c$:
\begin{enumerate}
  \item $\alpha(\mu_c)=0$ y $\omega(\mu_c)=\omega_0>0$ (un par cruza el eje imaginario; no es un cero doble);
  \item \emph{transversalidad:} $\alpha'(\mu_c)\neq0$;
  \item \emph{no degeneración:} el \clave{primer coeficiente de Lyapunov} $\ell_1$ (definido abajo) no es cero.
\end{enumerate}
Entonces de $\eq(\mu_c)$ nace una familia de órbitas periódicas, del lado de $\mu_c$ donde $\alpha'(\mu_c)(\mu-\mu_c)/\ell_1<0$, con amplitud proporcional a $\sqrt{\abs{\mu-\mu_c}}$ y período que tiende a $2\pi/\omega_0$ cuando $\mu\to\mu_c$. Si $\ell_1<0$ el ciclo es estable y la bifurcación es supercrítica; si $\ell_1>0$ es inestable y la bifurcación es subcrítica.
\end{teorema}

\notamargen{En la forma normal $\dot r=\alpha'\,(\mu-\mu_c)\,r+\ell_1r^3$, que es la del teorema hasta orden tres, el radio del ciclo es $r^2=-\alpha'(\mu-\mu_c)/\ell_1$. De ahí salen la ley de la raíz y la condición sobre el lado.}

Para calcular $\ell_1$, se lleva el sistema (con $\mu=\mu_c$ y el equilibrio en el origen) a la forma
\[
  \dot u=-\omega_0v+F(u,v),\qquad \dot v=\omega_0u+G(u,v),
\]
con $F$ y $G$ de orden dos o más. Entonces \citep[ec.~(3.4.11)]{GuHo83}
\begin{equation}
\begin{aligned}
  \ell_1&=\tfrac1{16}\bigl(F_{uuu}+F_{uvv}+G_{uuv}+G_{vvv}\bigr)\\
  &\quad+\tfrac1{16\omega_0}\bigl(F_{uv}(F_{uu}+F_{vv})-G_{uv}(G_{uu}+G_{vv})-F_{uu}G_{uu}+F_{vv}G_{vv}\bigr),
\end{aligned}
  \label{eq:bf-l1}
\end{equation}
con las derivadas evaluadas en el origen. Para la forma normal \eqref{eq:bf-hopfnormal}, $F=-u(u^2+v^2)$, $G=-v(u^2+v^2)$ y $\ell_1=\tfrac1{16}(-6-2-2-6)=-1$.

\begin{ejemplo}[label=ej:bf-l1circuito]{¿Supercrítica o subcrítica?}
Decida el tipo de la bifurcación del circuito \eqref{eq:bf-circuito} y prediga la amplitud del ciclo para $\mu$ pequeño.
\solucion En $\mu=0$: $\dot x_1=x_2-x_1^3$, $\dot x_2=-x_1$. Con $u=x_1$, $v=-x_2$: $\dot u=-v-u^3$, $\dot v=u$, que es la forma pedida con $\omega_0=1$, $F=-u^3$ y $G=0$. Solo sobrevive $F_{uuu}=-6$: $\ell_1=-\tfrac6{16}=-\tfrac38<0$, \alerta{supercrítica}. Con $\alpha'=\tfrac12$, el radio es $r^2\approx\tfrac12\mu/\tfrac38=\tfrac43\mu$: amplitud $2\sqrt{\mu/3}$. Con $\mu=0.1$ predice $0.3651$; la simulación da $0.3652$. Con $\mu=1$ predice $1.155$ contra $1.160$: la ley de la raíz es asintótica, pero aquí funciona bien incluso lejos. Es la amplitud que da el método de promedios para van der Pol.
\end{ejemplo}

\codigorepo{cap06_bifurcaciones/C6_Ej9_HopfSuperSub.m}{lyapunov}{El primer coeficiente de Lyapunov con derivadas numéricas}

En la práctica, sin fórmulas, la distinción se hace con una simulación: se toma $\mu$ un poco más allá del valor crítico y se arranca muy cerca del equilibrio. Si el estado se estaciona en un ciclo cuya amplitud crece como la raíz del exceso de $\mu$, la bifurcación es supercrítica; si se escapa, subcrítica. Para confirmar la subcrítica, se toma $\mu$ un poco antes del valor crítico y se integra \emph{hacia atrás} en el tiempo desde cerca del equilibrio: el ciclo inestable se vuelve atractor y aparece.

\begin{trampa}[label=tr:bf-centro]{``En $\mu=\mu_c$ la linealización es un centro: hay órbitas cerradas''}
\plausible{En $\mu=0$ la jacobiana de \eqref{eq:bf-hopfnormal} es $\bigl[\begin{smallmatrix}0&-1\\1&0\end{smallmatrix}\bigr]$, con valores propios $\pm j$. Es un centro, así que el sistema oscila alrededor del origen.}
Sin hiperbolicidad no hay Hartman--Grobman (\cref{tr:bi-centro}). Con $\mu=0$ la ecuación radial es $\dot r=-r^3$, que se integra exactamente: $r(t)=r_0/\sqrt{1+2r_0^2t}$. El origen es asintóticamente estable; desde $r_0=1$, $r(100)=1/\sqrt{201}=0.0705$. Pero el decaimiento es algebraico, como $1/\sqrt t$, no exponencial. No es ni un centro ni un foco exponencial; en la subcrítica, $\dot r=r^3$, es inestable.
\correcto{En el valor crítico deciden los términos cúbicos, es decir, $\ell_1$. Y cerca de él, la convergencia (o la divergencia) es lenta, lo que hace engañosas las simulaciones cortas.}
\end{trampa}

\begin{pausa}
Con $\mu=0$ en la forma normal supercrítica, ¿certifica $V=\tfrac12(x_1^2+x_2^2)$ la estabilidad asintótica del origen? ¿Certifica estabilidad exponencial?
\respuesta{$\dot V=r\dot r=-r^4=-(x_1^2+x_2^2)^2$, que es definida negativa: por el teorema de Lyapunov del \cref{cap:lyapunov}, el origen es asintóticamente estable, incluso globalmente porque $V$ es radialmente no acotada. Pero no exponencialmente: para eso haría falta $\dot V\le-cV$ cerca del origen, y $-r^4\le-\tfrac c2r^2$ falla para $r$ pequeño. Es consistente con $r(t)\sim1/\sqrt{2t}$.}
\end{pausa}

\begin{trampa}[label=tr:bf-periodo]{``El período del ciclo es $2\pi/\omega_0$''}
\plausible{El teorema de Hopf dice que el período es $2\pi/\omega_0$; en el circuito, $\omega_0=1$ y el período debería ser $2\pi$.}
El teorema dice que el período \emph{tiende} a $2\pi/\omega_0$ cuando $\mu\to\mu_c$. Lejos de $\mu_c$ cambia: en el circuito vale $6.287$ con $\mu=0.1$ y $6.663$ con $\mu=1$ (\cref{fig:bf-circuito}). En otros sistemas puede ser menor: un período medido de $5.9$ contra $2\pi/\omega_0=6.28$ no indica necesariamente un error de simulación.
\correcto{$T\to2\pi/\omega_0$ y la amplitud $\propto\sqrt{\abs{\mu-\mu_c}}$ son leyes asintóticas, válidas cerca de $\mu_c$. Se verifican en una gráfica log--log, no con un solo valor.}
\end{trampa}

Una versión del teorema que no necesita $\ell_1$, pero sí que la forma polar sea conocida, es la siguiente.

\begin{proposicion}[label=prop:bf-radial]{Ciclos de una ecuación radial}
Sea $\dot x_1=\mu x_1+x_2-x_1\,g(r)$, $\dot x_2=-x_1+\mu x_2-x_2\,g(r)$, con $r=\sqrt{x_1^2+x_2^2}$, $g$ de clase $\Cont1$, $g(0)=0$ y $g'(r)>0$ para $r>0$. Entonces: para $\mu<0$ el origen es un foco estable; para $\mu=0$ es asintóticamente estable, pero no exponencialmente; para $0<\mu<\sup g$ hay un único ciclo límite, el círculo $r=g^{-1}(\mu)$, que es estable y cuyo radio crece desde $0$ con $\mu$.
\end{proposicion}

\begin{demostracion}
En polares, $\dot r=r(\mu-g(r))$ y $\dot\theta=-1$. Para $\mu\le0$, $\mu-g(r)<0$ si $r>0$. Para $\mu>0$, $\mu-g(r)$ es positiva para $r<g^{-1}(\mu)$ y negativa después, porque $g$ es creciente: los círculos $r=c$ con $c<g^{-1}(\mu)$ se cruzan hacia afuera y los de radio mayor hacia adentro.
\end{demostracion}

\notamargen{La versión anterior de estas notas pedía ``$f(r)>0$ para $\mu>0$''; lo que hace falta es que $g$ sea creciente en $r$, para que el ciclo sea único y crezca con $\mu$.}

\subsection{Histéresis: la Hopf subcrítica en un sistema real}

En la subcrítica, ¿adónde va el estado cuando se escapa? En un sistema físico, términos de orden superior detienen el crecimiento. El modelo mínimo es
\begin{equation}
  \dot r=r\bigl(\mu+r^2-r^4\bigr),\qquad \dot\theta=1 .
  \label{eq:bf-histeresis}
\end{equation}
Los ciclos son las raíces de $\mu=r^4-r^2$, y un ciclo $r^*$ es estable si $\frac{\dif}{\dif r}\bigl(\mu+r^2-r^4\bigr)=2r^*-4r^{*3}<0$, es decir, si $r^*>1/\sqrt2$. Hay \emph{dos} bifurcaciones (\cref{fig:bf-histeresis}): en $\mu=0$, una Hopf subcrítica (el ciclo inestable se contrae hasta el origen); en $\mu=-\tfrac14$, $r=1/\sqrt2$, el ciclo estable y el inestable chocan y desaparecen: una \clave{silla-nodo de ciclos}, que se estudia en la \cref{sec:bf-ciclos}. Entre las dos, $-\tfrac14<\mu<0$, el sistema es \alerta{biestable}: el origen y el ciclo grande ($r=0.9420$ con $\mu=-0.1$) son estables.

Si $\mu$ sube despacio desde $-0.4$, el sistema se queda en reposo hasta pasar $\mu=0$ y salta a un ciclo de amplitud finita ($r\approx1$). Para apagar la oscilación no basta con regresar a $\mu=0$: hay que bajar hasta $-\tfrac14$. Es el esquema del \emph{flameo} (flutter) aeroelástico de un ala, del bombeo (surge) en compresores y de ciertas oscilaciones en sistemas de potencia: el sistema salta a una oscilación grande y no vuelve solo.

\begin{figure}[ht]
\centering
\includegraphics[width=\linewidth]{cap06/f_histeresis}
\caption{\eqref{eq:bf-histeresis}. Izquierda: diagrama de bifurcación (ciclos estables en verde, inestables en coral, el origen en negro). Derecha: una rampa lenta de $\mu$, de $-0.4$ a $0.1$ y de vuelta. Al subir, el salto llega con retardo ($\mu\approx0.05$); al bajar, el ciclo se apaga en $\mu\approx-\tfrac14$.}
\label{fig:bf-histeresis}
\end{figure}

\codigorepo{cap06_bifurcaciones/C6_Ej10_Histeresis.m}{rampa}{Rampa lenta de $\mu$ en la Hopf subcrítica con término quíntico}

\video{V6_3_histeresis_ciclos}{Salto y regreso: la histéresis de ciclos}

\begin{trampa}[label=tr:bf-retardo]{``Con una rampa lenta, el salto ocurre exactamente en $\mu=\mu_c$''}
\plausible{Si $\mu$ cambia muy despacio, el sistema está siempre casi en equilibrio y debe dejar el origen justo cuando este se vuelve inestable.}
Mientras $\mu<0$ el estado se acerca muchísimo al origen ($r$ decae exponencialmente), y cuando $\mu$ pasa a positivo necesita un tiempo largo para crecer desde ese valor diminuto. En la simulación de la \cref{fig:bf-histeresis} (rampa de $1.7\times10^{-4}$ por unidad de tiempo, con un piso de ruido de $10^{-6}$), el salto ocurre en $\mu\approx0.048$, no en $0$. Es el \clave{retardo de la bifurcación}: cuanto menos ruido, mayor el retardo. Al bajar, en cambio, el ciclo grande sí se apaga cerca de $-\tfrac14$ (en $-0.258$).
\correcto{Con un parámetro que cambia en el tiempo, el diagrama de bifurcación (que es estático) da los valores límite de una rampa infinitamente lenta \emph{y} con perturbaciones; el retardo depende de la velocidad de la rampa y del nivel de ruido.}
\end{trampa}

\notamargen{El modelo de Rosenzweig--MacArthur del \cref{cap:orbitas} (\cref{sec:oc-lotka}) tiene una Hopf en $K_H=3$: es la paradoja del enriquecimiento. El modelo de Nowak (\cref{sec:oc-nowak}), otra. Con el \cref{thm:bf-hopf} se puede decidir de qué tipo son (\cref{prob:bf-rm}).}

%% ---------------------------------------------------------------------
\section{Bifurcaciones de órbitas periódicas}
\label{sec:bf-ciclos}
%% ---------------------------------------------------------------------

Los ciclos límite también nacen y mueren, y no siempre en un equilibrio. Ya vimos uno: en \eqref{eq:bf-histeresis}, en $\mu=-\tfrac14$, un ciclo estable y uno inestable chocan y desaparecen lejos de cualquier equilibrio. Esas bifurcaciones \alerta{no se ven en los valores propios de ningún jacobiano}: el origen sigue siendo un foco estable a ambos lados de $\mu=-\tfrac14$. Para estudiarlas se necesita el mapa de Poincaré de la \cref{sec:bf-poincare}. Aquí describimos las tres más comunes en el plano \citep[\S8.4]{Str15}.

\subsection{Silla-nodo de ciclos}

Es la versión para ciclos de la silla-nodo: dos ciclos, uno estable y otro inestable, se acercan, se funden en un ciclo semiestable y desaparecen. Un ejemplo con todo explícito es \citep[cap.~7]{Sas99}
\begin{equation}
\begin{aligned}
  \dot x_1&=-x_1\sin\mu-x_2\cos\mu+(1-x_1^2-x_2^2)^2\,(x_1\cos\mu-x_2\sin\mu),\\
  \dot x_2&=\phantom{-}x_1\cos\mu-x_2\sin\mu+(1-x_1^2-x_2^2)^2\,(x_1\sin\mu+x_2\cos\mu),
\end{aligned}
\label{eq:bf-pliegue}
\end{equation}
que en polares es
\[
  \dot r=r\bigl((1-r^2)^2\cos\mu-\sin\mu\bigr),\qquad \dot\theta=(1-r^2)^2\sin\mu+\cos\mu .
\]
Para $\mu$ pequeño, $\dot\theta\approx1$. Si $\mu<0$, $\dot r>0$ para todo $r>0$: el origen es inestable y todas las soluciones se van al infinito. En $\mu=0$, $\dot r=r(1-r^2)^2\ge0$ y aparece un ciclo semiestable $r=1$ (atrae desde adentro, repele hacia afuera). Si $0<\mu<\pi/4$, los ciclos son las raíces de $(1-r^2)^2=\tan\mu$:
\[
  r=\sqrt{1\pm\sqrt{\tan\mu}} ,
\]
que con $\mu=0.1$ son $0.8266$ (estable) y $1.1475$ (inestable). Recorrida hacia $\mu$ decreciente, el ciclo estable choca con el inestable en $\mu=0$ y desaparece, y el estado, que estaba oscilando con amplitud finita, \alerta{se va al infinito} (\cref{fig:bf-pliegue}). Es una bifurcación peligrosa, y de nuevo no la anuncia ningún valor propio: el origen es un foco inestable a ambos lados.

\begin{figure}[ht]
\centering
\includegraphics[width=\linewidth]{cap06/f_pliegue}
\caption{La silla-nodo de ciclos \eqref{eq:bf-pliegue}. Izquierda: $\dot r$ contra $r$ para $\mu=-0.1$, $0$ y $0.1$. Centro: retrato con $\mu=0.1$, el ciclo estable (verde) dentro del inestable (coral). Derecha: diagrama de bifurcación de los radios.}
\label{fig:bf-pliegue}
\end{figure}

\begin{trampa}[label=tr:bf-cielo]{``Esta es la catástrofe de cielo azul''}
\plausible{Las diapositivas y la versión anterior de las notas llaman así a \eqref{eq:bf-pliegue}: un ciclo aparece ``de la nada'' al aumentar $\mu$.}
Lo que ocurre en \eqref{eq:bf-pliegue} es una silla-nodo de ciclos (también llamada \emph{pliegue} de ciclos): los dos ciclos nacen juntos, con radio finito, de un ciclo semiestable. La \clave{catástrofe de cielo azul} es otra cosa, descubierta por Turaev y Shilnikov en 1995~\citep{TuSh95}: una órbita periódica cuyo período y longitud crecen sin cota al acercarse al valor crítico, y que desaparece sin chocar con nada. No puede ocurrir en el plano; necesita al menos tres dimensiones. Y la ``bifurcación de cielo azul'' de Abraham y Shaw es la silla-nodo de \emph{equilibrios} (\cref{sec:bf-sillanodo}).
\correcto{En el plano, un ciclo que aparece con amplitud finita al lado de otro es una silla-nodo de ciclos. Los nombres con ``cielo azul'' son tres cosas distintas; conviene evitarlos o decir cuál se usa.}
\end{trampa}

\subsection{Período infinito y homoclínica}

En las otras dos, el período del ciclo crece sin cota al acercarse al valor crítico.
\begin{itemize}
  \item \textbf{Silla-nodo en un ciclo invariante} (SNIC, o bifurcación de período infinito). Una silla-nodo de equilibrios ocurre \emph{sobre} el ciclo. Es el oscilador no uniforme $\dot\theta=\omega-a\sin\theta$ del \cref{cap:orbitas}: para $a<\omega$ la fase da vueltas con período $2\pi/\sqrt{\omega^2-a^2}$ (\eqref{eq:oc-periodo}; $10.06$ con $\omega=2$, $a=1.9$), en $a=\omega$ aparece un equilibrio semiestable sobre el círculo y la oscilación se detiene. El período diverge como $1/\sqrt{\omega-a}$: la firma del fantasma de la \cref{sec:bf-sillanodo}, ahora en cada vuelta. La amplitud, en cambio, no se anula.
  \item \textbf{Homoclínica.} El ciclo crece hasta tocar una silla y se convierte en una órbita que sale de la silla y vuelve a ella (una conexión entre sillas, como las que prohíbe el \cref{thm:bf-peixoto}). El período diverge como $\abs{\ln\abs{\mu-\mu_c}}$, más despacio que en la SNIC.
\end{itemize}

\begin{ideaclave}
Cómo se apaga una oscilación, según lo que se mida al acercarse a $\mu_c$:
\begin{center}
\begin{tabular}{@{}lll@{}}
\toprule
\textbf{bifurcación} & \textbf{amplitud} & \textbf{período}\\
\midrule
Hopf supercrítica & $\to0$ como $\sqrt{\abs{\mu-\mu_c}}$ & $\to2\pi/\omega_0$\\
silla-nodo de ciclos & salto desde un valor finito & finito\\
SNIC (período infinito) & finita & $\sim1/\sqrt{\abs{\mu-\mu_c}}$\\
homoclínica & finita & $\sim\abs{\ln\abs{\mu-\mu_c}}$\\
\bottomrule
\end{tabular}
\end{center}
\end{ideaclave}

\begin{pausa}
En una simulación, una oscilación desaparece al cambiar un parámetro. Su amplitud no tiende a cero pero su período crece mucho justo antes. ¿Qué bifurcación descarta inmediatamente? ¿Cómo distinguiría entre las dos que quedan?
\respuesta{Descarta la Hopf supercrítica (la amplitud iría a cero con período finito) y la silla-nodo de ciclos (el período queda finito). Quedan la SNIC y la homoclínica. Se distinguen por la ley de crecimiento del período: graficando $T$ contra $\ln\abs{\mu-\mu_c}$ (recta para la homoclínica) y $\log T$ contra $\log\abs{\mu-\mu_c}$ (pendiente $-\tfrac12$ para la SNIC). También por la forma del ciclo justo antes: en la SNIC aparece un punto donde el estado ``se detiene'' (el fantasma), en la homoclínica una esquina, la silla.}
\end{pausa}

%% ---------------------------------------------------------------------
\section{Mapas de Poincaré}
\label{sec:bf-poincare}
%% ---------------------------------------------------------------------

¿Cómo se decide si un ciclo límite es estable, o se detecta una silla-nodo de ciclos, si no hay coordenadas polares que desacoplen el sistema? La idea de Poincaré es convertir el problema de una órbita periódica en el de un \emph{punto fijo} de un mapa: se mira el flujo solo cuando cruza una curva.

Sea $\gamma$ una órbita periódica de $\dot x=f(x)$, $x\in\R^2$, y $p\in\gamma$. Sea $\Sigma$ un segmento de curva que pasa por $p$ y que el flujo cruza \alerta{transversalmente} (el campo no es tangente a $\Sigma$ en ningún punto). Una trayectoria que arranca en $q\in\Sigma$, cerca de $p$, da una vuelta cerca de $\gamma$ y vuelve a cruzar $\Sigma$ en un punto $P(q)$ (\cref{fig:bf-seccion}).

\begin{definicion}[label=def:bf-poincare]{Mapa de Poincaré}
Con $\gamma$, $p$ y $\Sigma$ como arriba, y $U\subset\Sigma$ una vecindad de $p$ en $\Sigma$, el \clave{mapa de Poincaré} (o de primer retorno) es $P:U\to\Sigma$, $P(q)=\flujo{\tau(q)}q$, donde $\tau(q)>0$ es el primer instante en que la solución que arranca en $q$ vuelve a cruzar $\Sigma$ en el mismo sentido. $\tau(q)$ es el \clave{tiempo de retorno}. La sucesión $q,\ P(q),\ P^2(q),\dots$ es una \clave{sucesión de Poincaré}.
\end{definicion}

\notamargen{Que el retorno exista y que $P$ sea suave cerca de $p$ se sigue de la dependencia continua de las soluciones respecto a la condición inicial (\cref{cap:preliminares}) y del teorema de la función implícita aplicado a la ecuación ``$\flujo tq\in\Sigma$''. Lejos de $p$ el retorno puede no existir.}

\begin{figure}[ht]
\centering
\includegraphics[width=0.95\linewidth]{cap06/f_seccion}
\caption{Izquierda: la sección $\Sigma$ y la sucesión de Poincaré $A_0,A_1,A_2,\dots$ que converge al punto $B$ donde el ciclo estable corta a $\Sigma$. Derecha: la misma sucesión vista como iteración del mapa $P$ (diagrama de telaraña): de cada punto se sube a la gráfica de $P$ y se pasa horizontalmente a la diagonal.}
\label{fig:bf-seccion}
\end{figure}

Tres hechos convierten el mapa en una herramienta:
\begin{enumerate}
  \item \emph{Puntos fijos son órbitas periódicas.} Si $P(q^*)=q^*$, la solución que arranca en $q^*$ vuelve a $q^*$ en el tiempo $\tau(q^*)$: es periódica. Y cada órbita periódica que cruza $\Sigma$ da un punto fijo.
  \item \emph{El mapa es un sistema en tiempo discreto}, $q_{k+1}=P(q_k)$, y su estabilidad se decide igual que la de un equilibrio, pero con el círculo unitario en lugar del eje imaginario: si $\abs{P'(q^*)}<1$ el punto fijo es asintóticamente estable; si $\abs{P'(q^*)}>1$, inestable; si $\abs{P'(q^*)}=1$, la linealización no decide. El número $P'(q^*)$ es el \clave{multiplicador} del ciclo.
  \item \emph{La estabilidad del punto fijo es la del ciclo.} Si $q^*$ es asintóticamente estable para $P$, el ciclo $\gamma$ atrae a las trayectorias cercanas (es orbitalmente asintóticamente estable, en el sentido de la \cref{tr:oc-orbital}), y si es inestable, $\gamma$ también lo es.
\end{enumerate}
En el plano, $\Sigma$ es unidimensional y $P'$ es un número; además $P$ es creciente (dos trayectorias no se cruzan), así que el multiplicador es positivo. Una silla-nodo de ciclos es una silla-nodo del mapa: dos puntos fijos que chocan con $P'(q^*)=1$.

\begin{ejemplo}[label=ej:bf-mapa1]{Un mapa que se calcula exactamente}
Para $\dot x_1=x_2-x_1(x_1^2+x_2^2-1)$, $\dot x_2=-x_1-x_2(x_1^2+x_2^2-1)$, con $\Sigma$ el semieje $x_1>0$, calcule el mapa de Poincaré, muestre que hay una sola órbita periódica y clasifíquela \citep[Ejemplo~8.7.1]{Str15}.
\solucion En polares, $\dot r=r(1-r^2)$ y $\dot\theta=-1$: toda trayectoria vuelve a $\Sigma$ exactamente después de $\tau=2\pi$. Separando variables,
\[
  \int_{r_0}^{r}\Bigl[\frac1s+\frac{s}{1-s^2}\Bigr]\dif s=t
  \quad\Longrightarrow\quad
  \frac1{r(t)^2}-1=\Bigl(\frac1{r_0^2}-1\Bigr)e^{-2t},
\]
y con $t=2\pi$,
\[
  P(r_0)=\Bigl[1+\Bigl(\frac1{r_0^2}-1\Bigr)e^{-4\pi}\Bigr]^{-1/2} .
\]
El único punto fijo positivo es $r_0=1$ (si $P(r_0)=r_0$, entonces $(r_0^{-2}-1)(e^{-4\pi}-1)=0$). Derivando, $P'(1)=e^{-4\pi}\approx3.49\times10^{-6}<1$: el círculo unitario es un ciclo límite estable, y la distancia al ciclo se reduce en un factor de casi trescientos mil por vuelta.
\end{ejemplo}

\begin{ejemplo}[label=ej:bf-mapa2]{Un ciclo que nace en una Hopf, visto con el mapa}
Para $\dot x_1=\mu x_1+x_2-x_1\sqrt{x_1^2+x_2^2}$, $\dot x_2=-x_1+\mu x_2-x_2\sqrt{x_1^2+x_2^2}$, con $\mu>0$ y $\Sigma$ el semieje $x_1>0$, obtenga $P$ y estudie la sucesión que arranca en $r_0<\mu$ \citep[\S13.2]{JoSm07}.
\solucion Es la \cref{prop:bf-radial} con $g(r)=r$: en polares, $\dot r=r(\mu-r)$ y $\dot\theta=-1$. La ecuación radial es logística, con solución $r(t)=\mu r_0/\bigl(r_0+(\mu-r_0)e^{-\mu t}\bigr)$, y el retorno ocurre en $2\pi$:
\[
  P(r_0)=\frac{\mu r_0}{r_0+(\mu-r_0)e^{-2\pi\mu}} ,\qquad P(\mu)=\mu,\qquad P'(\mu)=e^{-2\pi\mu} .
\]
Con $\mu=0.2$, $P'(\mu)=e^{-0.4\pi}=0.2846$: la sucesión de Poincaré se acerca a $r=\mu$ desde adentro, monótonamente, y la distancia se multiplica por $0.28$ en cada vuelta (\cref{fig:bf-seccion}, derecha). Cuando $\mu\to0^+$, $P'(\mu)\to1$: el ciclo que nace en la Hopf es apenas atractor. Es la versión discreta de la convergencia lenta de la \cref{tr:bf-centro}.
\end{ejemplo}

\codigorepo{cap06_bifurcaciones/C6_Ej11_MapaPoincare.m}{mapa}{El mapa del \cref{ej:bf-mapa2}: fórmula contra simulación, y la telaraña}

\subsection{El multiplicador sin resolver la ecuación}

En los dos ejemplos el mapa salió de una solución explícita. En general no hay tal cosa, pero en el plano el multiplicador se puede calcular sin conocer $P$.

\begin{teorema}[label=thm:bf-divergencia]{El multiplicador en el plano \citep[\S3.4, Teorema~2]{Per01}}
Sea $\gamma(t)$ una órbita periódica de período $T$ de $\dot x=f(x)$, $f\in\Cont1(\R^2)$. Entonces, para cualquier sección transversal,
\[
  P'(q^*)=\exp\Bigl(\int_0^T\nabla\cdot f\bigl(\gamma(t)\bigr)\,\dif t\Bigr),\qquad \nabla\cdot f=\frac{\partial f_1}{\partial x_1}+\frac{\partial f_2}{\partial x_2} .
\]
En particular, $\gamma$ es estable si $\int_0^T\nabla\cdot f\,\dif t<0$ e inestable si es $>0$.
\end{teorema}

\notamargen{La divergencia mide cómo cambia el área de una nube de condiciones iniciales (\cref{thm:oc-bendixson} usa la misma cantidad). A lo largo del ciclo, una dirección (la tangente) no se contrae ni se expande; lo que se acumula es el cambio en la otra, que es el multiplicador. Es un caso particular de la teoría de Floquet.}

En el \cref{ej:bf-mapa2}, $\nabla\cdot f=2\mu-3r$, que sobre el ciclo $r=\mu$ vale $-\mu$; integrando durante $T=2\pi$, $P'(\mu)=e^{-2\pi\mu}$, como antes. Lo valioso es que también funciona cuando solo se tiene el ciclo numérico.

\begin{ejemplo}[label=ej:bf-vdp]{El mapa de Poincaré de van der Pol}
Para el circuito \eqref{eq:bf-circuito} con $\mu=1$ (van der Pol) y $\Sigma=\{x_2=0,\ x_1>0\}$, calcule numéricamente el punto fijo y el multiplicador.
\solucion Sobre $\Sigma$, $\dot x_2=-x_1<0$: el flujo cruza transversalmente, hacia abajo. Se integra desde $(q,0)$ hasta el siguiente cruce de $x_2=0$ hacia abajo (con un evento de \texttt{ode45}), y eso define $P(q)$. Una búsqueda de raíces de $P(q)-q$ da $q^*=1.1081$ (no es el máximo de $x_1$ sobre el ciclo, que es $1.160$ y ocurre fuera de $\Sigma$), con tiempo de retorno $6.663$, el período. Con diferencias finitas, $P'(q^*)=8.60\times10^{-4}$, que coincide con $\exp\int_0^T(1-3x_1^2)\,\dif t$ calculado a lo largo del ciclo: el ciclo de van der Pol es muy atractor. Y el tiempo de retorno cambia con $q$: desde $q=0.5$ y desde $q=2.5$ no es el mismo.
\end{ejemplo}

\begin{figure}[ht]
\centering
\includegraphics[width=\linewidth]{cap06/f_vdp_poincare}
\caption{El mapa de Poincaré de van der Pol ($\mu=1$) en $\Sigma=\{x_2=0,\ x_1>0\}$. Izquierda: tres trayectorias y sus cruces con $\Sigma$. Centro: la gráfica de $P$ (casi plana cerca del punto fijo: $P'(q^*)\approx10^{-3}$) y la telaraña desde $q=0.2$. Derecha: el tiempo de retorno $\tau(q)$, que no es constante.}
\label{fig:bf-vdp}
\end{figure}

\codigorepo{cap06_bifurcaciones/C6_Ej12_PoincareVanDerPol.m}{vdp}{Mapa de Poincaré numérico de van der Pol y su multiplicador por dos caminos}

\video{V6_4_poincare}{La sucesión de Poincaré de van der Pol}

\begin{trampa}[label=tr:bf-retorno]{``El tiempo entre retornos es constante''}
\plausible{En los dos ejemplos el retorno ocurre siempre en $2\pi$, y como el sistema es invariante en el tiempo, el tiempo de una vuelta debería ser siempre el mismo.}
Que el sistema sea invariante en el tiempo quiere decir que el instante en que se \emph{empieza} no importa, no que todas las vueltas duren lo mismo. En los ejemplos el tiempo de retorno es constante porque $\dot\theta$ no depende de $r$. En van der Pol depende de $q$, y en general $\tau(q)$ es una función, igual al período solo en el punto fijo.
\correcto{$P(q)=\flujo{\tau(q)}q$ con $\tau$ que depende de $q$. Para calcular $P$ numéricamente se usa un evento que detecte el cruce de $\Sigma$, no un tiempo fijo.}
\end{trampa}

\begin{trampa}[label=tr:bf-uno]{``$P'(q^*)=1$: el ciclo es neutral, como un centro''}
\plausible{Si el multiplicador es $1$, las trayectorias vecinas no se acercan ni se alejan, así que el ciclo está rodeado de órbitas cerradas.}
$P'(q^*)=1$ es el caso en que la linealización del mapa no decide, igual que un valor propio nulo. En el ciclo semiestable de \eqref{eq:bf-pliegue} con $\mu=0$ el multiplicador es $1$ (sobre $r=1$, $\partial\dot r/\partial r=0$), y el ciclo atrae desde adentro y repele hacia afuera. Y en el \cref{ej:bf-mapa2}, $P'(\mu)\to1$ cuando el ciclo nace. Deciden los términos de orden superior de $P$.
\correcto{$\abs{P'(q^*)}=1$ marca un ciclo no hiperbólico: el candidato a una bifurcación de ciclos. Para clasificarlo hay que mirar $P''$ (o más).}
\end{trampa}

Los mapas de Poincaré son aún más útiles en sistemas forzados periódicamente, $\dot x=f(t,x)$ con $f(t+T,x)=f(t,x)$: allí el plano de fase no sirve (las trayectorias se cruzan, porque el campo cambia con el tiempo), pero el \clave{mapa estroboscópico} $x(kT)\mapsto x\bigl((k+1)T\bigr)$ es un mapa invariante en el tiempo. Una solución periódica del sistema forzado es un punto fijo del mapa, y el caos del Duffing forzado (\cref{tr:bi-caos}) se ve como una nube de puntos con estructura fractal. Otra técnica para estudiar la estabilidad de órbitas periódicas, la de Floquet, linealiza el sistema a lo largo del ciclo y estudia la matriz de transición resultante \citep[cap.~7]{Sas99}; en el plano, el \cref{thm:bf-divergencia} es su resultado principal.

\begin{pausa}
En la \cref{fig:bf-seccion}, si se hubiera escogido otra sección $\Sigma'$, el punto fijo sería otro, $B'$. ¿Cambia el multiplicador?
\respuesta{No. El punto fijo depende de la sección (es donde $\Sigma'$ corta al ciclo), pero el multiplicador no: el \cref{thm:bf-divergencia} lo expresa con una integral sobre el ciclo completo, que no depende de dónde se corte. Más en general, los mapas de dos secciones distintas son conjugados por el flujo entre ellas, y los mapas conjugados tienen la misma derivada en el punto fijo.}
\end{pausa}

%% ---------------------------------------------------------------------
\section{Para llevar}
\label{sec:bf-cierre}
%% ---------------------------------------------------------------------

\begin{parallevar}
\begin{enumerate}
  \item \textbf{Qué se rompe:} una bifurcación local es una pérdida de hiperbolicidad. Si la jacobiana es invertible, el equilibrio se continúa suavemente con $\mu$ y conserva su tipo. En $\R^2$ hay dos rutas: $\delta\to0$ (silla-nodo, transcrítica, horquilla) y $\tau\to0$ con $\delta>0$ (Hopf).
  \item \textbf{Receta:} equilibrios, jacobiano, cruce, \emph{transversalidad} y \emph{no degeneración}. Sin los dos últimos pasos no se sabe qué bifurcación hay.
  \item \textbf{Silla-nodo:} $\dot x=\mu\pm x^2$; dos equilibrios chocan y desaparecen. Es genérica. Del lado sin equilibrios queda un fantasma que demora las soluciones un tiempo $\sim1/\sqrt{\abs{\mu-\mu_c}}$; del otro, la recuperación se vuelve lenta.
  \item \textbf{Transcrítica y horquilla:} intercambio de estabilidad y aparición de pares simétricos. No son genéricas: las protegen un equilibrio que existe para todo $\mu$ o una simetría. Supercrítica es segura; subcrítica es peligrosa y, con términos de orden superior, produce histéresis.
  \item \textbf{En el plano,} la Ruta~1 ocurre sobre la variedad central, donde vale el catálogo de la recta; hay que reducir sobre $W^c$, no sobre un eje.
  \item \textbf{Hopf:} un par complejo cruza el eje imaginario con $\alpha'\neq0$ y nace un ciclo de amplitud $\propto\sqrt{\abs{\mu-\mu_c}}$ y período $\to2\pi/\omega_0$. El signo de $\ell_1$ decide: negativo, supercrítica (ciclo estable); positivo, subcrítica (ciclo inestable, salto e histéresis).
  \item \textbf{Ciclos:} los ciclos también nacen y mueren lejos de los equilibrios (silla-nodo de ciclos, SNIC, homoclínica), sin que lo anuncie ningún valor propio.
  \item \textbf{Mapa de Poincaré:} órbitas periódicas son puntos fijos de $P$; el ciclo es estable si $\abs{P'(q^*)}<1$. En el plano, $P'(q^*)=\exp\int_0^T\nabla\cdot f\,\dif t$, y no depende de la sección.
\end{enumerate}
\end{parallevar}

\begin{trampa}[title={\faBug\hspace{0.45em}Las trampas del capítulo, en una caja}]{}
\begin{itemize}
  \item $\det J=0$ es necesario, no suficiente (\cref{tr:bf-det}); con $\delta=0$ la traza no decide la estabilidad (\cref{tr:bf-traza}).
  \item Sin equilibrio no significa sin demora: el fantasma (\cref{tr:bf-fantasma}).
  \item La subcrítica no lleva a las ramas $\pm\sqrt{-\mu}$ (\cref{tr:bf-subcritica}); transcrítica y horquilla no son genéricas (\cref{tr:bf-generica}).
  \item Se reduce sobre la variedad central, no poniendo $x_2=0$ (\cref{tr:bf-proyectar}).
  \item Una silla-nodo con otro atractor no es reversible: histéresis (\cref{tr:bf-reversible}); con una rampa, el salto llega tarde (\cref{tr:bf-retardo}).
  \item En el valor crítico de Hopf no hay un centro (\cref{tr:bf-centro}); el período es $2\pi/\omega_0$ solo en el límite (\cref{tr:bf-periodo}).
  \item La silla-nodo de ciclos no es la catástrofe de cielo azul (\cref{tr:bf-cielo}).
  \item El tiempo de retorno no es constante (\cref{tr:bf-retorno}); $P'=1$ no es un centro (\cref{tr:bf-uno}).
\end{itemize}
\end{trampa}

\begin{profundiza}[Correcciones respecto a la versión anterior (Handout5 y las diapositivas de bifurcaciones)]
(i) Horquilla subcrítica en 2-D: los equilibrios para $\mu<0$ son $(\pm\sqrt{-\mu},0)$, no $(\pm\sqrt\mu,0)$, y la viñeta ``en $\mu=0$ \dots{} aparecen dos nuevos nodos estables'' estaba copiada del caso supercrítico (\cref{sec:bf-plano}). (ii) ``$\eq(\mu)$ es suave siempre y cuando la jacobiana no tenga valores propios nulos'': es una condición suficiente (función implícita), no necesaria, y un valor propio nulo no hace por sí solo que una bifurcación sea peligrosa (\cref{prop:bf-implicita}, \cref{tr:bf-det}). (iii) Interruptor genético: el origen es un nodo estable para todo $a,b>0$, también si $a=b$ (valor propio doble, pero estable); el umbral entre los dos estados es la variedad estable de la silla; las tasas son de \emph{degradación}. (iv) El teorema de Hopf estaba enunciado solo para la forma polar y con la hipótesis ``$f(r)>0$ para $\mu>0$''; se reemplazó por la \cref{prop:bf-radial} ($g$ creciente) y se agregó el enunciado completo con transversalidad y el coeficiente $\ell_1$ (\cref{thm:bf-hopf}). (v) El ejemplo llamado ``Hopf subcrítica'' y repetido como ``Blue Sky Catastrophe'' no es una Hopf (el origen es inestable a ambos lados) sino una silla-nodo de ciclos, y la catástrofe de cielo azul es otra cosa (\cref{tr:bf-cielo}). Además, en las ecuaciones cartesianas faltaba el cuadrado: con $(1-x_1^2-x_2^2)$ en lugar de $(1-x_1^2-x_2^2)^2$ la ecuación radial es $\dot r=r((1-r^2)\cos\mu-\sin\mu)$, con un solo ciclo. (vi) La horquilla subcrítica con término quíntico: ``existe un punto $\mu_s$ en el cual al aumentar $x$ el sistema se vuelve estable'' se reemplazó por las sillas-nodo exactas en $\mu=-\tfrac14$ y el ciclo de histéresis. (vii) Se unificó el parámetro ($\mu$; la versión anterior mezclaba $r$ y $\mu$ en la horquilla). (viii) Mapas de Poincaré: ``el tiempo entre retornos es constante'' es falso en general (\cref{tr:bf-retorno}); ``son más útiles para sistemas no autónomos'' se precisó como el mapa estroboscópico de los sistemas forzados periódicamente; se resolvió el ejemplo que estaba pendiente (\cref{ej:bf-mapa2}). (ix) Estabilidad estructural: el teorema citado requiere que el flujo cruce transversalmente la frontera de $D$, se debe a Andronov y Pontryagin en el disco y a Peixoto en superficies compactas, y la hiperbolicidad incluye a las órbitas periódicas (\cref{thm:bf-peixoto}). (x) Atribuciones: la bifurcación de Hopf es de Eberhard Hopf (1942), con antecedentes de Poincaré y Andronov; el teorema del índice es de Heinz Hopf. ``Poicaré'' por Poincaré. (xi) Se agregaron: la proposición de la función implícita, la receta de cinco pasos, las condiciones de la silla-nodo, el fantasma y la lentitud crítica, las bifurcaciones imperfectas, la horquilla escondida, el coeficiente de Lyapunov, el retardo con rampa lenta, las bifurcaciones de período infinito y homoclínica, el teorema de la divergencia para el multiplicador y el mapa de Poincaré numérico de van der Pol.
\end{profundiza}

%% ---------------------------------------------------------------------
\section{Problemas}
\label{sec:bf-problemas}
%% ---------------------------------------------------------------------

Los problemas marcados con \etiquetaM{} requieren \matlab/\simulink; los marcados con \etiquetaEstrella{} son de profundización. Cada uno tiene una trampa: búsquela antes de dar la respuesta.

\subsection*{En la recta}

\begin{problema}[ponerle nombre]\label{prob:bf-clasificar}
Encuentre los valores de bifurcación, dibuje el diagrama y diga de qué tipo es cada bifurcación: (a) $\dot x=1+\mu x+x^2$; (b) $\dot x=\mu x-\ln(1+x)$, con $x>-1$; (c) $\dot x=\mu x-\dfrac{x}{1+x^2}$. En (c), compare con la forma normal de la horquilla \emph{antes} de calcular.
\end{problema}

\begin{problema}[cuánto dura el colapso]\label{prob:bf-colapso}
Para la cosecha $\dot x=x(1-x)-h$: (a) con $h=0.26$, calcule el tiempo que tarda la población en pasar de $x=1$ a $x=0$; (b) repita con $h=0.3$; (c) \etiquetaM{} grafique ese tiempo contra $h-\tfrac14$ en escala log--log. ¿Qué pendiente espera? ¿Por qué una cosecha ``apenas'' excesiva es la más engañosa?
\end{problema}

\begin{problema}[una columna imperfecta]\label{prob:bf-columna}
Un modelo de la columna con una imperfección pequeña es $\dot x=\mu x-x^3+h$, $h>0$. (a) Muestre que para $\mu<3(h/2)^{2/3}$ hay un solo equilibrio. (b) ¿Hacia qué lado se pandea la columna al aumentar $\mu$ lentamente desde $\mu=-1$? (c) ¿Puede la columna terminar pandeada hacia el otro lado? ¿Cómo?
\end{problema}

\subsection*{En el plano}

\begin{problema}[otro interruptor]\label{prob:bf-griffith}
En el interruptor genético \eqref{eq:bf-griffith} con $b=0.5$: (a) halle $a_c$; (b) con $a=0.8$, halle los tres equilibrios y clasifíquelos; (c) un estudiante dice que con $a=0.8$ el gen siempre termina encendido porque ``el nodo alto es estable''. ¿Qué le falta?
\end{problema}

\begin{problema}[Hopf a mano]\label{prob:bf-hopfmano}
Para $\dot x_1=\mu x_1-x_2+ax_1^3$, $\dot x_2=x_1+\mu x_2$: (a) verifique las hipótesis 1 y 2 del \cref{thm:bf-hopf} en $\mu=0$; (b) calcule $\ell_1$ en función de $a$; (c) ¿para qué $a$ es supercrítica? Prediga el radio del ciclo; (d) ¿qué pasa con $a=0$?
\end{problema}

\begin{problema}[la paradoja del enriquecimiento, otra vez]\label{prob:bf-rm}
En el modelo de Rosenzweig--MacArthur \eqref{eq:oc-rm} con $\beta=2$, $\gamma=1$: (a) muestre que en el equilibrio de coexistencia $\tau=(K-3)/(2K)$ y $\delta=(K-1)/(2K)$; (b) verifique la transversalidad en $K=3$ y calcule $\omega_0$ y el período del ciclo al nacer; (c) \etiquetaM{} con \cmd{primerLyapunov} decida si la Hopf es supercrítica o subcrítica. ¿Es consistente con la simulación del \cref{cap:orbitas}?
\end{problema}

\begin{problema}[dónde se apaga la oscilación]\label{prob:bf-quintica}
Para $\dot r=r(\mu+cr^2-r^4)$, $\dot\theta=1$, con $c>0$: (a) encuentre la silla-nodo de ciclos; (b) ¿en qué intervalo de $\mu$ es biestable el sistema? (c) ¿Qué pasa con ese intervalo cuando $c\to0^+$? ¿Y si $c<0$?
\end{problema}

\subsection*{Mapas de Poincaré}

\begin{problema}[dos ciclos]\label{prob:bf-dosciclos}
Para $\dot r=-r(1-r)(2-r)$, $\dot\theta=1$: (a) encuentre los ciclos; (b) muestre que, para una ecuación $\dot r=g(r)$, $\dot\theta=1$, el multiplicador de un ciclo $r^*$ es $e^{2\pi g'(r^*)}$, y compárelo con el \cref{thm:bf-divergencia}; (c) clasifique los ciclos. ¿Qué pasa con las trayectorias que arrancan entre ellos?
\end{problema}

\begin{problema}[el circuito con el mapa]\label{prob:bf-mapacircuito}
\etiquetaM{} Adapte \cmd{C6\_Ej12\_PoincareVanDerPol} al circuito \eqref{eq:bf-circuito} con $\mu=0.1$ y $\mu=0.5$. (a) Encuentre el punto fijo y el multiplicador. (b) Compárelos con lo que predice la forma normal de Hopf, $P'\approx e^{-2\pi\mu}$. ¿Para cuál $\mu$ es mejor la predicción? ¿Por qué?
\end{problema}

\begin{problema}[el retardo del interruptor]\label{prob:bf-retardo}
\etiquetaM{} Con \cmd{crearModelo\_InterruptorGenetico.m} ($b=0.35$), haga subir $a$ desde $1$ con velocidad $\varepsilon=10^{-3}$, $1.25\times10^{-4}$ y $1.56\times10^{-5}$, y mida en qué valor de $a$ cae $x$ por debajo de $0.7$. (a) ¿Se apaga el gen en $a_c=1.4286$? (b) ¿Cómo cambia el retardo $a-a_c$ cuando $\varepsilon$ se divide por $8$?
\end{problema}

\begin{problema}[sin duplicación de período en el plano]\label{prob:bf-plano}
\etiquetaEstrella{} (a) Muestre que el mapa de Poincaré de un sistema plano es creciente. (b) Deduzca que el multiplicador de un ciclo es positivo y que $P$ no puede tener puntos de período dos ($P^2(q)=q$ con $P(q)\neq q$). (c) ¿Qué bifurcación de ciclos que existe en $\R^3$ queda descartada en el plano?
\end{problema}

%% ---------------------------------------------------------------------
\section{Soluciones seleccionadas}
%% ---------------------------------------------------------------------

\paragraph{Problema~\ref{prob:bf-clasificar}}
(a) $f=0$ tiene raíz doble cuando $\mu^2=4$: hay sillas-nodo en $\mu=2$ ($\eq=-1$) y $\mu=-2$ ($\eq=1$); hay dos equilibrios si $\abs\mu>2$ y ninguno si $\abs\mu<2$. (b) $x=0$ es equilibrio para todo $\mu$, con $f_x(0)=\mu-1$: transcrítica en $\mu=1$; la otra rama, $\mu x=\ln(1+x)$, pasa por el origen con $\eq\approx2(1-\mu)$. (c) $f$ es impar y $x=0$ es equilibrio: el tipo es de horquilla, en $\mu=1$ ($f_x(0)=\mu-1$). Los equilibrios no nulos cumplen $1+x^2=1/\mu$, así que existen para $0<\mu<1$, del lado en que el origen es estable, y allí $f_x=2\mu(1-\mu)>0$: son inestables. Es una horquilla \emph{subcrítica}, aunque $-x/(1+x^2)$ ``se parece'' a $-x+x^3$: la trampa es que el signo de $x^3$ en la expansión, $+x^3$, es el que decide.

\paragraph{Problema~\ref{prob:bf-colapso}}
(a) $x(1-x)-0.26=-\bigl(0.01+(x-\tfrac12)^2\bigr)$, y $T=\int_0^1\dif x/\bigl(0.01+(x-\tfrac12)^2\bigr)=20\arctan5=27.47$. (b) Con $a=\sqrt{0.05}$, $T=(2/a)\arctan(0.5/a)=10.29$. (c) Pendiente $-\tfrac12$ cerca de $h=\tfrac14$. Con una cosecha apenas excesiva, la población se queda muchos años cerca de $x=\tfrac12$ (el fantasma) y parece sostenible, hasta que colapsa.

\paragraph{Problema~\ref{prob:bf-columna}}
(a) La cúbica $\mu x-x^3+h$ tiene tres raíces reales si y solo si su discriminante $4\mu^3-27h^2$ es positivo, es decir, $\mu>3(h/2)^{2/3}$. (b) Hacia $x>0$: el único equilibrio, para $\mu$ pequeño, es positivo y se continúa sin bifurcar (es la rama continua de la \cref{fig:bf-imperfectas}). (c) Sí, pero no por el cambio lento de $\mu$: la rama negativa nace en una silla-nodo aislada en $\mu=3(h/2)^{2/3}$, y para llegar a ella hay que empujar la columna a través de la variedad estable de la silla.

\paragraph{Problema~\ref{prob:bf-griffith}}
(a) $a_c=1/(2b)=1$. (b) $\eq_\pm=(1\pm\sqrt{1-0.16})/0.8$, es decir, $2$ y $0.5$, además del origen. $\tau=-1.3$; en $x=2$, $\delta=2\cdot3/25=0.24>0$ y $\tau^2-4\delta=0.73>0$: nodo estable; en $x=0.5$, $\delta=-0.24$: silla; el origen, nodo estable. (c) Le falta el otro atractor: el origen también es estable, y el estado final depende de qué lado de la variedad estable de la silla esté la condición inicial. Un gen que empieza apagado se queda apagado.

\paragraph{Problema~\ref{prob:bf-hopfmano}}
(a) $J=\bigl[\begin{smallmatrix}\mu&-1\\1&\mu\end{smallmatrix}\bigr]$, $\lambda=\mu\pm j$: $\alpha(0)=0$, $\omega_0=1$, $\alpha'=1$. (b) Ya está en la forma pedida con $F=au^3$, $G=0$: $\ell_1=6a/16=3a/8$. (c) Supercrítica si $a<0$, con $r^2\approx-\mu/\ell_1=8\mu/(3\abs a)$. (d) Con $a=0$ el sistema es lineal: en $\mu=0$ hay una familia de órbitas cerradas (un centro) y para $\mu\neq0$ ninguna. No hay ciclo límite: la hipótesis de no degeneración falla, y el teorema no se puede aplicar.

\paragraph{Problema~\ref{prob:bf-rm}}
(a) Con $v^*=1$ y $p^*=2-2/K$, la jacobiana es $\bigl[\begin{smallmatrix}(K-3)/(2K)&-1/2\\(K-1)/K&0\end{smallmatrix}\bigr]$. (b) $\dif\tau/\dif K=1/6$ en $K=3$, así que $\alpha'=1/12\neq0$; $\omega_0=\sqrt\delta=1/\sqrt3$ y el período al nacer es $2\pi\sqrt3=10.88$. (c) $\ell_1<0$: supercrítica. Con las coordenadas que usa \cmd{primerLyapunov}, $\ell_1=-1/24\approx-0.042$; la magnitud depende de la escala de las coordenadas, el signo no. El ciclo existe para $K>3$, es estable y crece desde cero, como en la simulación con $K=5$ del \cref{cap:orbitas}.

\paragraph{Problema~\ref{prob:bf-quintica}}
(a) Los ciclos cumplen $\mu=r^4-cr^2$, que tiene su mínimo en $r^2=c/2$: silla-nodo de ciclos en $\mu=-c^2/4$. (b) $-c^2/4<\mu<0$. (c) Cuando $c\to0^+$ el intervalo se cierra y la Hopf pasa de subcrítica a supercrítica: es una bifurcación de codimensión dos (Bautin), donde $\ell_1=0$. Con $c<0$ la Hopf es supercrítica y no hay biestabilidad.

\paragraph{Problema~\ref{prob:bf-dosciclos}}
(a) $r=1$ y $r=2$. (b) Linealizando $\dot r=g(r)$ alrededor de $r^*$, $\delta r(2\pi)=e^{2\pi g'(r^*)}\delta r(0)$. La divergencia en polares es $g'(r)+g(r)/r$, que sobre el ciclo vale $g'(r^*)$; integrada durante $2\pi$ da lo mismo. (c) $g(r)=-2r+3r^2-r^3$: $g'(1)=1$ (multiplicador $e^{2\pi}$, inestable) y $g'(2)=-2$ (multiplicador $e^{-4\pi}$, estable). Entre los dos ciclos $\dot r>0$: las trayectorias se alejan del interior y se acercan al exterior.

\paragraph{Problema~\ref{prob:bf-mapacircuito}}
Con $\mu=0.1$: $q^*=0.3650$ y $P'=0.5331$, contra $e^{-0.2\pi}=0.5335$. Con $\mu=0.5$: $q^*=0.8068$ y $P'=0.0392$, contra $0.0432$. La predicción es mejor con $\mu=0.1$ porque la forma normal es una aproximación de orden tres válida cerca del valor crítico.

\paragraph{Problema~\ref{prob:bf-retardo}}
(a) No: con $\varepsilon=10^{-3}$ el gen cae en $a\approx a_c+0.059$; con $1.25\times10^{-4}$, en $a_c+0.017$; con $1.56\times10^{-5}$, en $a_c+0.0045$. El estado se demora en el fantasma de la silla-nodo mientras $a$ sigue subiendo. (b) Se divide por algo menos de $4$ ($3.5$ y $3.7$), y la razón se acerca a $4=8^{2/3}$: el retardo de una silla-nodo con rampa lenta escala como $\varepsilon^{2/3}$.

\paragraph{Problema~\ref{prob:bf-plano}}
(a) Si $q_1<q_2$ en $\Sigma$, las trayectorias que arrancan en ellos no se pueden cruzar, así que sus siguientes cruces con $\Sigma$ conservan el orden. (b) $P$ creciente implica $P'\ge0$, y como $P'(q^*)$ es una exponencial, $P'>0$. Si $P(q)>q$, entonces $P^2(q)>P(q)>q$ (y al revés): las sucesiones de Poincaré son monótonas y no pueden volver a $q$ en dos pasos. (c) La duplicación de período, que necesita un multiplicador que pase por $-1$. Por eso en el plano no hay cascadas de duplicación ni caos (Poincaré--Bendixson).

\respuestaspausas

